\documentclass[UTF8,a4paper,11pt]{ctexart}
\usepackage{geometry}
\geometry{left=1.8cm,right=1.8cm,top=1.8cm,bottom=1.8cm}
\usepackage{amsmath,amssymb,mathtools,bm}
\usepackage{hyperref}
\usepackage{xcolor}
\usepackage{enumitem}
\usepackage{fancyhdr}
\usepackage{fvextra}
\usepackage{longtable,array,booktabs}
\usepackage{titlesec}
\usepackage{setspace}
\setstretch{1.08}
\setlength{\parindent}{2em}
\setlength{\parskip}{0.35em}
\setcounter{tocdepth}{4}
\pagestyle{fancy}
\setlength{\headheight}{14pt}
\fancyhf{}
\fancyhead[L]{boys band cry}
\fancyhead[R]{\thepage}
\fancyfoot[C]{}
\definecolor{codebg}{RGB}{248,248,248}
\RecustomVerbatimEnvironment{Verbatim}{Verbatim}{fontsize=\small,frame=single,rulecolor=\color{black},formatcom=\color{black},baselinestretch=1}
\title{boys band cry\\\large template}
\author{}
\date{}
\begin{document}
\maketitle
\thispagestyle{fancy}
\renewcommand{\contentsname}{目录}
\setcounter{tocdepth}{3}
\tableofcontents
\clearpage
\section*{1. 数论}\addcontentsline{toc}{section}{1. 数论}

\subsection*{1.1 素数}\addcontentsline{toc}{subsection}{1.1 素数}

\subsubsection*{1.1.1 线性筛素数}\addcontentsline{toc}{subsubsection}{1.1.1 线性筛素数}

\begin{Verbatim}[fontsize=\small,frame=single,breaklines=true,breakanywhere=true]
// bool vis[N]
void getprime(int MAXN){
    for(int i=2;i<=MAXN;i++){
        if(!vis[i]) prim[++prim[0]]=i;
        for(int j=1;j<=prim[0]&&prim[j]*i<=MAXN;++j){
            vis[prim[j]*i]=1;
            if(i%prim[j]==0) break;
        }
    }
}
// int vis[N]
void getprim(int MAXN){
    for(int i=2;i<=MAXN;++i){
        if(!vis[i]) vis[i]=prim[++prim[0]]=i;
        for(int j=1;j<=prim[0]&&i*prim[j]<=MAXN;++j){
            vis[i*prim[j]]=prim[j];
        }
    }
}
\end{Verbatim}
\subsubsection*{1.1.2 埃筛素数}\addcontentsline{toc}{subsubsection}{1.1.2 埃筛素数}

\begin{Verbatim}[fontsize=\small,frame=single,breaklines=true,breakanywhere=true]
void getprim(int MAXN){
    int cnt(0);
    for(int i=2;i<=MAXN;++i){
        if(!vis[i]){
            prim[++cnt]=i;
            for(int j=i+i;j<=MAXN;j+=i) vis[j]=1;
        }
    }
}
\end{Verbatim}

\subsubsection*{1.1.3.1 素性检验（试除法）}\addcontentsline{toc}{subsubsection}{1.1.3.1 素性检验（试除法）}

\begin{Verbatim}[fontsize=\small,frame=single,breaklines=true,breakanywhere=true]
bool isPrime(int a){
    if(a<2) return 0;
  for(int i=2;i<a;++i)
    if(a%i==0) return 0;
  return 1;
}
\end{Verbatim}
\subsubsection*{1.1.3.2 素性检验（Miller–Rabin）}\addcontentsline{toc}{subsubsection}{1.1.3.2 素性检验（Miller–Rabin）}

二次探测定理：若 $p$ 为奇素数，则 $x^2\equiv1 \pmod{p}$ 或者 $x^2\equiv p-1 \pmod{p}$.

\begin{Verbatim}[fontsize=\small,frame=single,breaklines=true,breakanywhere=true]
bool MillerRabin(int n){
    int test_time(30);
    if(n<3||n%2==0) return n==2;
    if(n%3==0) return n==3;
    int u(n-1),t(0);
    for(;u%2==0;u>>=1,++t);
    for(;test_time--;){
        int a(rand()%(n-3)+2);
        int v(qpow(a,u,n));// 快速幂
        if(v==1) continue;
        int s;
        for(s=0;s<t;++s){
            if(v==n-1) break;
            v=(long long) v*v%n;
        }
        if(s==t) return 0;
    }
    return 1;
}
\end{Verbatim}
\subsubsection*{1.1.4 质因数分解}\addcontentsline{toc}{subsubsection}{1.1.4 质因数分解}

\begin{Verbatim}[fontsize=\small,frame=single,breaklines=true,breakanywhere=true]
vector<int> BreakDown(int MAXN){
    vector<int> res;
    for(int i=2;i*i<=MAXN;++i){
        if(MAXN%i==0){
            for(;MAXN%i==0;MAXN/=i);
            res.push_back(i);
        }
    }
    if(MAXN>1) res.push_back(MAXN);
    return res;
}
\end{Verbatim}
\subsection*{1.2 欧拉函数}\addcontentsline{toc}{subsection}{1.2 欧拉函数}

\subsubsection*{1.2.1 根号欧拉函数}\addcontentsline{toc}{subsubsection}{1.2.1 根号欧拉函数}

设 $n=\prod^s_{i=1}p_i^{k_i}$，则有：

\[
\begin{aligned}
\varphi(n)&=\prod^s_{i=1}\varphi(p_i^{k_i})\\
    &=\prod^s_{i=1}(p_i-1)p_i^{k_i-1}\\
    &=\prod^s_{i=1}(1-\frac{1}{p_i})p_i^{k_i}\\
    &=n\prod^s_{i=1}(1-\frac{1}{p_i})
\end{aligned}
\]

\begin{Verbatim}[fontsize=\small,frame=single,breaklines=true,breakanywhere=true]
int euler_phi(int n){
    int ans(n);
    for(int i=2;i*i<=n;i++)
        if(n%i==0){
            ans=ans/i*(i-1);
            for(n%i==0) n/=i;
        }
    if(n>1) ans=ans/n*(n-1);
    return ans;
}
\end{Verbatim}
推论：对任意不全为 $0$ 的整数 $m,n,\varphi(mn)\varphi(\gcd(m,n))=\varphi(m)\varphi(n)\gcd(m,n)$.

\subsubsection*{1.2.2 欧拉求和}\addcontentsline{toc}{subsubsection}{1.2.2 欧拉求和}

\[
n=\sum_{d|n}\varphi(d)
\]

\subsubsection*{1.2.2 欧拉求和与反演}\addcontentsline{toc}{subsubsection}{1.2.2 欧拉求和与反演}

\[
\begin{aligned}
n&=\sum_{d|n}\varphi(d)\\
\gcd(a,b)&=\sum_{d|\gcd(a,b)}\varphi(d)\\
\gcd(a,b)&=\sum_{d}[d|a][d|b]\varphi(d)\\
\sum_{i=1}^n \gcd(i,n)&=\sum_{d|n}\lfloor\frac{n}{d}\rfloor\varphi(d)\\
\end{aligned}
\]

常用莫比乌斯反演：若 $f(n)=\sum_{d|n}g(d)$，则

\[
g(n)=\sum_{d|n}\mu(d)f\left(\frac{n}{d}\right)
\]

以及

\[
\sum_{d|n}\mu(d)=[n=1]
\]

\subsection*{1.3 筛法}\addcontentsline{toc}{subsection}{1.3 筛法}

\subsubsection*{1.3.1 筛法求欧拉函数}\addcontentsline{toc}{subsubsection}{1.3.1 筛法求欧拉函数}

\begin{Verbatim}[fontsize=\small,frame=single,breaklines=true,breakanywhere=true]
void pre_eular(int MAXN){
    phi[1]=1;
    for(int i=2;i<=MAXN;++i){
        if(!phi[i]) prim[++prim[0]]=i,phi[i]=i-1;
        for(int j=1;j<=prim[0]&&i*prim[j]<=MAXN;++j){
            if(!(i%prim[j])){
                phi[i*prim[j]]=prim[j]*phi[i];
                break;
            }
            phi[i*prim[j]]=(prim[j]-1)*phi[i];
        }
    }
}
\end{Verbatim}
\subsubsection*{1.3.2 筛法求莫比乌斯函数}\addcontentsline{toc}{subsubsection}{1.3.2 筛法求莫比乌斯函数}

\begin{Verbatim}[fontsize=\small,frame=single,breaklines=true,breakanywhere=true]
void pre_mu(int MAXN){
    mu[1]=1;
    for(int i=2;i<=MAXN;i++){
        if(!vis[i]) prime[++cnt]=i,mu[i]=-1;
        for(int j=1;j<=cnt&&i*prime[j]<=MAXN;j++){
            vis[i*prime[j]]=1;
            if(i%prime[j]==0) break;
            mu[i*prime[j]]=-mu[i];
        }
    }
}
\end{Verbatim}
\subsubsection*{1.3.3 筛法求约数个数函数}\addcontentsline{toc}{subsubsection}{1.3.3 筛法求约数个数函数}

\begin{Verbatim}[fontsize=\small,frame=single,breaklines=true,breakanywhere=true]
void pre_d(int MAXN){
    d[1]=1;
    for(int i=2;i<=MAXN;++i){
        if(!vis[i]) prim[++prim[0]]=i,d[i]=2,num[i]=1;
        for(int j=1;j<=prim[0]&&i*prim[j]<=MAXN;++j){
            vis[prim[j]*i]=1;
            if(i%prim[j]==0){
                num[i*prim[j]]=num[i]+1;
                d[i*prim[j]]=d[i]/(num[i]+1)*(num[i]+2);
                break;
            }
            num[i*prim[j]]=1;
            d[i*prim[j]]=d[i]<<1;
        }
    }
}
\end{Verbatim}

\subsubsection*{1.3.4 筛法求约数和}\addcontentsline{toc}{subsubsection}{1.3.4 筛法求约数和}
生成函数乘积，取质因子随机组合。
\begin{Verbatim}[fontsize=\small,frame=single,breaklines=true,breakanywhere=true]
void pre_f(int MAXN){
    // g[i]为i的最小质因子p的k次定比数列和
    g[1]=f[1]=1;
    for(int i=2;i<=MAXN;++i){
        if(!vis[i]) prim[++prim[0]]=i,g[i]=f[i]=i+1;
        for(int j=1;j<=prim[0]&&i*prim[j]<=MAXN;++j){
            vis[prim[j]*i]=1;
            if(i%prim[j]==0){
                g[i*prim[j]]=g[i]*prim[j]+1;
                f[i*prim[j]]=f[i]/g[i]*g[i*prim[j]];
                break;
            }
            g[i*prim[j]]=1+prim[j];
            f[i*prim[j]]=f[i]*(prim[j]+1);
        }
    }
}
\end{Verbatim}
\subsection*{1.4 最大公约数}\addcontentsline{toc}{subsection}{1.4 最大公约数}

\subsubsection*{1.4.1 欧几里得算法}\addcontentsline{toc}{subsubsection}{1.4.1 欧几里得算法}

\begin{Verbatim}[fontsize=\small,frame=single,breaklines=true,breakanywhere=true]
int gcd(int a,int b){return !b?a:gcd(b,a%b);}
int gcd(int a,int b){
  for(int tmp;b;)
    tmp=a,a=b,b=tmp%b;
  return a;
}
\end{Verbatim}

\subsubsection*{1.4.2 扩展欧几里得算法}\addcontentsline{toc}{subsubsection}{1.4.2 扩展欧几里得算法}

\begin{Verbatim}[fontsize=\small,frame=single,breaklines=true,breakanywhere=true]
void Exgcd(ll a,ll b,ll &x,ll &y){
    if(!b) x=1,y=0;
    else Exgcd(b,a%b,y,x),y-=a/b*x;
}
\end{Verbatim}

\subsection*{1.5 模算术类}\addcontentsline{toc}{subsection}{1.5 模算术类}

\subsubsection*{1.5.1 快速幂}\addcontentsline{toc}{subsubsection}{1.5.1 快速幂}

\begin{Verbatim}[fontsize=\small,frame=single,breaklines=true,breakanywhere=true]
ll qpow(ll a,ll b,ll p){
    ll res(1);
    for(;b;b>>=1,a=a*a%p) if(b&1) res=res*a%p;
    return res%p;
}
\end{Verbatim}

\subsubsection*{1.5.2 逆元}\addcontentsline{toc}{subsubsection}{1.5.2 逆元}

\paragraph*{1.5.2.1 费马小定理}\addcontentsline{toc}{paragraph}{1.5.2.1 费马小定理}
\begin{Verbatim}[fontsize=\small,frame=single,breaklines=true,breakanywhere=true]
int inv=qpow(a,p-2,p);// p为质数
\end{Verbatim}
\paragraph*{1.5.2.2 线性求逆元}\addcontentsline{toc}{paragraph}{1.5.2.2 线性求逆元}
\[
i^{-1}\equiv-\lfloor\frac{p}{i}\rfloor\times(p\text{ mod } i)^{-1} \pmod{p}
\]

\begin{Verbatim}[fontsize=\small,frame=single,breaklines=true,breakanywhere=true]
inv[1]=1;
for(int i=2;i<p;++i)
    inv[i]=(p-p/i)*inv[p%i]%p;
\end{Verbatim}

\paragraph*{1.5.2.2 扩欧求逆元}\addcontentsline{toc}{paragraph}{1.5.2.2 扩欧求逆元}

\begin{Verbatim}[fontsize=\small,frame=single,breaklines=true,breakanywhere=true]
void Exgcd(ll a,ll b,ll &x,ll &y){
    if(!b) x=1,y=0;
    else Exgcd(b,a%b,y,x),y-=a/b*x;
}
int main(){
    ll x,y;
    Exgcd(a,p,x,y);
    x=(x%p+p)%p;//x是a在mod p下的逆元
}
\end{Verbatim}
\paragraph*{1.5.2.3 阶乘逆元}\addcontentsline{toc}{paragraph}{1.5.2.3 阶乘逆元}

\[
\frac{1}{i!}\times i=\frac{1}{(i-1)!}
\]

得到：

\[
\frac{1}{i!}\times (i-1)!=\frac{1}{i}
\]

更一般的将 $i$ 替换为 $a_i$ 也可线性求列表逆元。

\subsubsection*{1.5.3 卢卡斯定理}\addcontentsline{toc}{subsubsection}{1.5.3 卢卡斯定理}

\begin{Verbatim}[fontsize=\small,frame=single,breaklines=true,breakanywhere=true]
ll Lucas(ll n,ll k,ll p) {
    if(!k) return 1ll;
  return(C(n%p,k%p,p)*Lucas(n/p,k/p,p))%p;
}// p为质数
\end{Verbatim}
\subsubsection*{1.5.4 扩展卢卡斯定理}\addcontentsline{toc}{subsubsection}{1.5.4 扩展卢卡斯定理}

\subsection*{1.6 中国剩余定理}\addcontentsline{toc}{subsection}{1.6 中国剩余定理}

\subsubsection*{1.6.1 CRT（模数互质）}\addcontentsline{toc}{subsubsection}{1.6.1 CRT（模数互质）}

若 $m_i$ 两两互质，求

\[
x\equiv a_i\pmod{m_i}
\]

则唯一解模 $M=\prod m_i$。

\begin{Verbatim}[fontsize=\small,frame=single,breaklines=true,breakanywhere=true]
ll CRT(vector<ll> a,vector<ll> m){
    ll M=1,ans=0;
    for(auto x:m) M*=x;
    for(int i=0;i<(int)a.size();++i){
        ll Mi=M/m[i],x,y;
        Exgcd(Mi,m[i],x,y);
        x=(x%m[i]+m[i])%m[i];
        ans=(ans+a[i]*Mi%M*x)%M;
    }
    return (ans%M+M)%M;
}
\end{Verbatim}

\subsubsection*{1.6.2 exCRT（模数不互质）}\addcontentsline{toc}{subsubsection}{1.6.2 exCRT（模数不互质）}

逐个合并：

\[
\begin{cases}
x\equiv a_1\pmod{m_1}\\
x\equiv a_2\pmod{m_2}
\end{cases}
\]

若 $a_2-a_1$ 不能被 $\gcd(m_1,m_2)$ 整除，则无解。

\begin{Verbatim}[fontsize=\small,frame=single,breaklines=true,breakanywhere=true]
bool exCRT(ll a1,ll m1,ll a2,ll m2,ll &a,ll &m){
    ll x,y;
    ll g=gcd(m1,m2);
    if((a2-a1)%g) return 0;
    Exgcd(m1,m2,x,y);
    ll p=m2/g;
    x=(x%p+p)%p;
    ll k=((a2-a1)/g%p+p)%p*x%p;
    a=a1+m1*k;
    m=m1/g*m2;
    a=(a%m+m)%m;
    return 1;
}
bool merge(const vector<ll> &a,const vector<ll> &m,ll &a_,ll &m_){
    if(a.empty()) return 0;
    a_=a[0],m_=m[0];
    for(int i=1;i<(int)a.size();++i)
        if(!exCRT(a)) return 0;
    return 1;
}
\end{Verbatim}

\subsection*{1.7 数论分块}\addcontentsline{toc}{subsection}{1.7 数论分块}

数论分块常用于处理 $\sum_{i=1}^n\lfloor\frac{n}{i}\rfloor$、$\sum f(\lfloor n/i\rfloor)$ 之类的式子。

\begin{Verbatim}[fontsize=\small,frame=single,breaklines=true,breakanywhere=true]
for(ll l=1,r;l<=n;l=r+1){
    r=n/(n/l);
    // [l,r] 内 n/i 的值相同
}
\end{Verbatim}

\subsection*{1.8 杜教筛}\addcontentsline{toc}{subsection}{1.8 杜教筛}

预处理后可在 $O(n^{2/3})$ 左右的复杂度内求 $\sum_{i=1}^n\mu(i)$ 和 $\sum_{i=1}^n\varphi(i)$。调用 \texttt{init(limit)} 后使用 \texttt{sumMu(n)}、\texttt{sumPhi(n)}。

\VerbatimInput[fontsize=\small,frame=single,breaklines=true,breakanywhere=true]{数学/数论/杜教筛.cpp}

\subsection*{1.9 Min\_25 筛}\addcontentsline{toc}{subsection}{1.9 Min\_25 筛}

以下模板以 $\sum_{i=1}^n\varphi(i)$ 为例，结果模 $10^9+7$；只需调用 \texttt{sumPhi(n)}。

\VerbatimInput[fontsize=\small,frame=single,breaklines=true,breakanywhere=true]{数学/数论/Min_25筛.cpp}

\subsection*{1.10 洲阁筛}\addcontentsline{toc}{subsection}{1.10 洲阁筛}

洲阁筛保留所有 $\lfloor n/i\rfloor$ 状态。以下以 $\sum_{i=1}^n d(i)$ 为例，结果模 $10^9+7$；调用 \texttt{sumDivisorCount(n)}。

\VerbatimInput[fontsize=\small,frame=single,breaklines=true,breakanywhere=true]{数学/数论/洲阁筛.cpp}


\section*{2. 线性代数}\addcontentsline{toc}{section}{2. 线性代数}

\subsection*{2.1 高斯消元}\addcontentsline{toc}{subsection}{2.1 高斯消元}

\begin{Verbatim}[fontsize=\small,frame=single,breaklines=true,breakanywhere=true]
int Gauss(vector<vector<double>> a,vector<double> &ans){
    int n=a.size(),m=a[0].size()-1,r=0;
    for(int c=0;c<m&&r<n;++c){
        int p=r;
        for(int i=r;i<n;++i) if(fabs(a[i][c])>fabs(a[p][c])) p=i;
        if(fabs(a[p][c])<1e-9) continue;
        swap(a[p],a[r]);
        double div=a[r][c];
        for(int j=c;j<=m;++j) a[r][j]/=div;
        for(int i=0;i<n;++i) if(i!=r&&fabs(a[i][c])>1e-9){
            double t=a[i][c];
            for(int j=c;j<=m;++j) a[i][j]-=t*a[r][j];
        }
        ++r;
    }
    ans.assign(m,0);
    for(int i=0;i<r;++i){
        int pos=0;
        while(pos<m&&fabs(a[i][pos])<1e-9) ++pos;
        if(pos<m) ans[pos]=a[i][m];
    }
    return r;
}
\end{Verbatim}

\subsection*{2.2 矩阵快速幂}\addcontentsline{toc}{subsection}{2.2 矩阵快速幂}

\begin{Verbatim}[fontsize=\small,frame=single,breaklines=true,breakanywhere=true]
struct Mat{
    int n;
    vector<vector<ll>> a;
    Mat(int _n=0,int id=0):n(_n),a(_n,vector<ll>(_n,0)){
        if(id) for(int i=0;i<n;++i) a[i][i]=1;
    }
};

Mat operator*(const Mat &A,const Mat &B){
    Mat C(A.n);
    for(int i=0;i<A.n;++i)
        for(int k=0;k<A.n;++k) if(A.a[i][k])
            for(int j=0;j<A.n;++j)
                C.a[i][j]+=A.a[i][k]*B.a[k][j];
    return C;
}
\end{Verbatim}

\subsection*{2.3 线性基}\addcontentsline{toc}{subsection}{2.3 线性基}

\begin{Verbatim}[fontsize=\small,frame=single,breaklines=true,breakanywhere=true]
ll basis[61];
void insert(ll x){
    for(int i=60;i>=0;--i){
        if(!(x>>i&1)) continue;
        if(!basis[i]) return basis[i]=x,void();
        x^=basis[i];
    }
}
ll queryMax(ll x){
    for(int i=60;i>=0;--i) x=max(x,x^basis[i]);
    return x;
}
\end{Verbatim}

\subsection*{2.4 行列式}\addcontentsline{toc}{subsection}{2.4 行列式}

\begin{Verbatim}[fontsize=\small,frame=single,breaklines=true,breakanywhere=true]
long double det(vector<vector<long double>> a){
    int n=a.size();
    long double ans=1;
    for(int i=0;i<n;++i){
        int p=i;
        for(int j=i;j<n;++j) if(fabsl(a[j][i])>fabsl(a[p][i])) p=j;
        if(fabsl(a[p][i])<1e-18) return 0;
        if(i!=p){ swap(a[i],a[p]); ans=-ans; }
        ans*=a[i][i];
        for(int j=i+1;j<n;++j){
            long double t=a[j][i]/a[i][i];
            for(int k=i;k<n;++k) a[j][k]-=t*a[i][k];
        }
    }
    return ans;
}
\end{Verbatim}

\subsection*{2.5 矩阵求逆}\addcontentsline{toc}{subsection}{2.5 矩阵求逆}

\begin{Verbatim}[fontsize=\small,frame=single,breaklines=true,breakanywhere=true]
bool inverse(vector<vector<double>> a,vector<vector<double>> &inv){
    int n=a.size();
    inv.assign(n,vector<double>(n,0));
    for(int i=0;i<n;++i) inv[i][i]=1;
    for(int c=0;c<n;++c){
        int p=c;
        for(int i=c;i<n;++i) if(fabs(a[i][c])>fabs(a[p][c])) p=i;
        if(fabs(a[p][c])<1e-9) return 0;
        swap(a[p],a[c]);
        swap(inv[p],inv[c]);
        double div=a[c][c];
        for(int j=0;j<n;++j) a[c][j]/=div,inv[c][j]/=div;
        for(int i=0;i<n;++i) if(i!=c){
            double t=a[i][c];
            for(int j=0;j<n;++j) a[i][j]-=t*a[c][j],inv[i][j]-=t*inv[c][j];
        }
    }
    return 1;
}
\end{Verbatim}

\subsection*{2.6 递推转矩阵}\addcontentsline{toc}{subsection}{2.6 递推转矩阵}

线性递推 $f_n=c_1f_{n-1}+\cdots+c_kf_{n-k}$ 可写成矩阵快速幂。

\begin{Verbatim}[fontsize=\small,frame=single,breaklines=true,breakanywhere=true]
// 形如 f[n] = sum c[i]*f[n-i]
// 构造 k*k 转移矩阵后快速幂
\end{Verbatim}


\section*{3. 组合数学}\addcontentsline{toc}{section}{3. 组合数学}

\subsection*{3.1 组合数预处理}\addcontentsline{toc}{subsection}{3.1 组合数预处理}

\begin{Verbatim}[fontsize=\small,frame=single,breaklines=true,breakanywhere=true]
ll C[2005][2005];
void initC(int n){
    for(int i=0;i<=n;++i){
        C[i][0]=C[i][i]=1;
        for(int j=1;j<i;++j)
            C[i][j]=(C[i-1][j-1]+C[i-1][j])%p;
    }
}
\end{Verbatim}

\subsection*{3.2 组合数取模（阶乘）}\addcontentsline{toc}{subsection}{3.2 组合数取模（阶乘）}

\begin{Verbatim}[fontsize=\small,frame=single,breaklines=true,breakanywhere=true]
ll fac[N],ifac[N];
void initFac(int n){
    fac[0]=1;
    for(int i=1;i<=n;++i) fac[i]=fac[i-1]*i%p;
    ifac[n]=qpow(fac[n],p-2,p);
    for(int i=n;i>=1;--i) ifac[i-1]=ifac[i]*i%p;
}

ll C(ll n,ll m){
    if(n<m||m<0) return 0;
    return fac[n]*ifac[m]%p*ifac[n-m]%p;
}
\end{Verbatim}

\subsection*{3.3 容斥原理}\addcontentsline{toc}{subsection}{3.3 容斥原理}

\[
\left|\bigcup_i A_i\right|=\sum|A_i|-\sum|A_i\cap A_j|+\sum|A_i\cap A_j\cap A_k|-\cdots
\]

常用于“至少一个满足”“没有任何满足”的计数问题。

\subsection*{3.4 Catalan 数}\addcontentsline{toc}{subsection}{3.4 Catalan 数}

\[
Cat_n=\frac{1}{n+1}\binom{2n}{n}
\]

\begin{Verbatim}[fontsize=\small,frame=single,breaklines=true,breakanywhere=true]
ll Catalan(int n){
    return C(2*n,n)-C(2*n,n-1);
}
\end{Verbatim}
\subsection*{3.5 排列与错排}\addcontentsline{toc}{subsection}{3.5 排列与错排}

\begin{Verbatim}[fontsize=\small,frame=single,breaklines=true,breakanywhere=true]
ll der[N];
void initDer(int n){
    der[0]=1; der[1]=0;
    for(int i=2;i<=n;++i) der[i]=(i-1)*(der[i-1]+der[i-2])%p;
}
\end{Verbatim}

错排公式：

\[
!n=n!\sum_{i=0}^{n}\frac{(-1)^i}{i!}
\]

\subsection*{3.6 Stirling 数}\addcontentsline{toc}{subsection}{3.6 Stirling 数}

第二类 Stirling 数 $S(n,k)$ 表示把 $n$ 个不同元素划分为 $k$ 个非空集合的方案数。

\[
S(n,k)=S(n-1,k-1)+kS(n-1,k)
\]

\begin{Verbatim}[fontsize=\small,frame=single,breaklines=true,breakanywhere=true]
ll S[205][205];
void initS(int n){
    S[0][0]=1;
    for(int i=1;i<=n;++i)
        for(int j=1;j<=i;++j)
            S[i][j]=(S[i-1][j-1]+j*S[i-1][j])%p;
}
\end{Verbatim}

\subsection*{3.7 Bell 数}\addcontentsline{toc}{subsection}{3.7 Bell 数}

Bell 数表示集合划分总数：

\[
B_n=\sum_{k=0}^n S(n,k)
\]

\subsection*{3.8 多重集与插板法}\addcontentsline{toc}{subsection}{3.8 多重集与插板法}

将 $n$ 个相同物品分到 $k$ 个盒子：

\[
\binom{n+k-1}{k-1}
\]

若每盒至少 1 个：

\[
\binom{n-1}{k-1}
\]

\subsection*{3.9 Burnside 引理}\addcontentsline{toc}{subsection}{3.9 Burnside 引理}

设群 $G$ 作用在集合 $X$ 上，则不同本质状态数为

\[
\frac{1}{|G|}\sum_{g\in G}|\mathrm{Fix}(g)|
\]

\subsection*{3.10 Pólya 定理}\addcontentsline{toc}{subsection}{3.10 Pólya 定理}

用于“旋转/翻转等价”的颜色计数，可配合循环指标生成函数使用。


\section*{4. 多项式与生成函数}\addcontentsline{toc}{section}{4. 多项式与生成函数}

\subsection*{4.1 NTT}\addcontentsline{toc}{subsection}{4.1 NTT}

\begin{Verbatim}[fontsize=\small,frame=single,breaklines=true,breakanywhere=true]
void ntt(vector<ll> &a,int inv){
    int n=a.size();
    // 先做 bit-reverse 置换
    for(int len=1;len<n;len<<=1){
        ll wn=qpow(G,(p-1)/(len<<1),p);
        if(inv==-1) wn=qpow(wn,p-2,p);
        for(int i=0;i<n;i+=len<<1){
            ll w=1;
            for(int j=0;j<len;++j){
                ll x=a[i+j],y=w*a[i+j+len]%p;
                a[i+j]=(x+y)%p;
                a[i+j+len]=(x-y+p)%p;
                w=w*wn%p;
            }
        }
    }
    if(inv==-1){
        ll invn=qpow(n,p-2,p);
        for(auto &x:a) x=x*invn%p;
    }
}
\end{Verbatim}
\subsection*{4.1.1 NTT半家桶（ln/exp/pow/sqrt）}\addcontentsline{toc}{subsection}{4.1.1 NTT半家桶（ln/exp/pow/sqrt）}

\begin{Verbatim}[fontsize=\small,frame=single,breaklines=true,breakanywhere=true]
// 1) 模数 p 为 NTT 友好质数（如 998244353），原根 G
// 2) qpow(a,b,p)
// 3) ntt(vector<ll>& a, int inv)（inv=1 正变换，-1 逆变换）

vector<ll> polyCut(vector<ll> a,int n){
    if((int)a.size()>n) a.resize(n);
    return a;
}

vector<ll> polyDeriv(const vector<ll> &a){
    if(a.empty()) return {};
    vector<ll> b(max(0,(int)a.size()-1));
    for(int i=1;i<(int)a.size();++i) b[i-1]=a[i]*i%p;
    return b;
}

vector<ll> polyInteg(const vector<ll> &a){
    vector<ll> b((int)a.size()+1);
    for(int i=0;i<(int)a.size();++i) b[i+1]=a[i]*qpow(i+1,p-2,p)%p;
    return b;
}

vector<ll> polyMul(vector<ll> a,vector<ll> b){
    if(a.empty()||b.empty()) return {};
    int n=(int)a.size(),m=(int)b.size(),lim=1;
    while(lim<n+m-1) lim<<=1;
    a.resize(lim); b.resize(lim);
    ntt(a,1); ntt(b,1);
    for(int i=0;i<lim;++i) a[i]=a[i]*b[i]%p;
    ntt(a,-1);
    a.resize(n+m-1);
    return a;
}

// A[0] != 0
vector<ll> polyInv(const vector<ll> &A,int n){
    vector<ll> B(1,qpow(A[0],p-2,p));
    for(int len=2;len<=n;len<<=1){
        vector<ll> F(min((int)A.size(),len));
        for(int i=0;i<(int)F.size();++i) F[i]=A[i];
        vector<ll> T=polyMul(polyMul(B,B),F);
        B.resize(len);
        for(int i=0;i<len;++i){
            ll v=(i<(int)T.size()?T[i]:0);
            B[i]=(2LL*B[i]%p-v+p)%p;
        }
    }
    B.resize(n);
    return B;
}

// ln(A) = integral(A' / A), 要求 A[0] = 1
vector<ll> polyLn(const vector<ll> &A,int n){
    vector<ll> d=polyDeriv(A);
    vector<ll> invA=polyInv(A,n);
    vector<ll> q=polyMul(d,invA);
    q=polyCut(q,n-1);
    vector<ll> res=polyInteg(q);
    return polyCut(res,n);
}

// exp(A), 要求 A[0] = 0
vector<ll> polyExp(const vector<ll> &A,int n){
    vector<ll> B(1,1); // exp(0)=1
    for(int len=2;len<=n;len<<=1){
        vector<ll> lnB=polyLn(B,len);
        vector<ll> C(len);
        for(int i=0;i<len;++i){
            ll ai=(i<(int)A.size()?A[i]:0);
            ll li=(i<(int)lnB.size()?lnB[i]:0);
            C[i]=(ai-li+p)%p;
        }
        C[0]=(C[0]+1)%p;
        B=polyMul(B,C);
        B.resize(len);
    }
    B.resize(n);
    return B;
}

// sqrt(A)，要求给定 s0^2 ≡ A[0] (mod p)，且 A[0]!=0
vector<ll> polySqrt(const vector<ll> &A,int n,ll s0){
    ll inv2=(p+1)/2;
    vector<ll> B(1,s0%p);
    for(int len=2;len<=n;len<<=1){
        vector<ll> invB=polyInv(B,len);
        vector<ll> F(min((int)A.size(),len));
        for(int i=0;i<(int)F.size();++i) F[i]=A[i];
        vector<ll> T=polyMul(F,invB);
        B.resize(len);
        for(int i=0;i<len;++i){
            ll ti=(i<(int)T.size()?T[i]:0);
            B[i]=(B[i]+ti)%p*inv2%p;
        }
    }
    B.resize(n);
    return B;
}

// pow(A,k)：返回 A^k mod x^n
// 要点：处理前导零。若 A = x^t * B, 则 A^k = x^(t*k) * B^k
vector<ll> polyPow(vector<ll> A,ll k,int n){
    if(n==0) return {};
    int t=0;
    while(t<(int)A.size()&&A[t]==0) ++t;
    if(t==(int)A.size()) return vector<ll>(n,0);
    if(1LL*t*k>=n) return vector<ll>(n,0);

    ll lead=A[t];
    ll invLead=qpow(lead,p-2,p);
    vector<ll> B(A.begin()+t,A.end());
    for(auto &x:B) x=x*invLead%p;

    int m=n-(int)(1LL*t*k);
    B=polyCut(B,m);

    // B[0] = 1
    vector<ll> lnB=polyLn(B,m);
    for(auto &x:lnB) x=x*(k%(p-1))%p;
    vector<ll> C=polyExp(lnB,m);

    ll coef=qpow(lead,k%(p-1),p);
    for(auto &x:C) x=x*coef%p;

    vector<ll> res(n,0);
    int shift=(int)(1LL*t*k);
    for(int i=0;i<(int)C.size()&&i+shift<n;++i) res[i+shift]=C[i];
    return res;
}
\end{Verbatim}

\subsection*{4.2 多项式求逆}\addcontentsline{toc}{subsection}{4.2 多项式求逆}

设 $B(x)=A(x)^{-1}$，则

\[
B(x)\cdot A(x)\equiv 1\pmod{x^n}
\]

常用递推：$B_{k+1}=B_k(2-A B_k)\bmod x^{2^k}$。

\subsection*{4.3 多项式求导与积分}\addcontentsline{toc}{subsection}{4.3 多项式求导与积分}

\begin{Verbatim}[fontsize=\small,frame=single,breaklines=true,breakanywhere=true]
vector<ll> deriv(vector<ll> a){
    vector<ll> b(max(0,(int)a.size()-1));
    for(int i=1;i<(int)a.size();++i) b[i-1]=a[i]*i%p;
    return b;
}

vector<ll> integ(vector<ll> a){
    vector<ll> b(a.size()+1);
    for(int i=0;i<(int)a.size();++i) b[i+1]=a[i]*qpow(i+1,p-2,p)%p;
    return b;
}
\end{Verbatim}

\subsection*{4.4 生成函数常见结论}\addcontentsline{toc}{subsection}{4.4 生成函数常见结论}

\begin{itemize}[leftmargin=2em,itemsep=0.2em,topsep=0.2em]
\item 普通生成函数：$F(x)=\sum a_n x^n$
\item 指数生成函数：$F(x)=\sum a_n\frac{x^n}{n!}$
\item 卷积对应乘法：若 $c_n=\sum_{i+j=n}a_ib_j$，则 $C(x)=A(x)B(x)$

\end{itemize}
\subsection*{4.5 FWT（按位卷积）}\addcontentsline{toc}{subsection}{4.5 FWT（按位卷积）}

\subsubsection*{4.5.1 异或卷积}\addcontentsline{toc}{subsubsection}{4.5.1 异或卷积}

\begin{Verbatim}[fontsize=\small,frame=single,breaklines=true,breakanywhere=true]
void fwt_xor(vector<ll> &a,int inv){
    int n=a.size();
    for(int len=1;len<n;len<<=1)
        for(int i=0;i<n;i+=len<<1)
            for(int j=0;j<len;++j){
                ll x=a[i+j],y=a[i+j+len];
                a[i+j]=(x+y)%p;
                a[i+j+len]=(x-y+p)%p;
            }
    if(inv==-1){
        ll invn=qpow(n,p-2,p);
        for(auto &x:a) x=x*invn%p;
    }
}
\end{Verbatim}

\subsubsection*{4.5.2 或卷积}\addcontentsline{toc}{subsubsection}{4.5.2 或卷积}

\begin{Verbatim}[fontsize=\small,frame=single,breaklines=true,breakanywhere=true]
void fwt_or(vector<ll> &a,int inv){
    int n=a.size();
    for(int len=1;len<n;len<<=1)
        for(int i=0;i<n;i+=len<<1)
            for(int j=0;j<len;++j)
                if(inv==1) a[i+j+len]=(a[i+j+len]+a[i+j])%p;
                else a[i+j+len]=(a[i+j+len]-a[i+j]+p)%p;
}
\end{Verbatim}

\subsubsection*{4.5.3 与卷积}\addcontentsline{toc}{subsubsection}{4.5.3 与卷积}

\begin{Verbatim}[fontsize=\small,frame=single,breaklines=true,breakanywhere=true]
void fwt_and(vector<ll> &a,int inv){
    int n=a.size();
    for(int len=1;len<n;len<<=1)
        for(int i=0;i<n;i+=len<<1)
            for(int j=0;j<len;++j)
                if(inv==1) a[i+j]=(a[i+j]+a[i+j+len])%p;
                else a[i+j]=(a[i+j]-a[i+j+len]+p)%p;
}
\end{Verbatim}

\subsection*{4.6 拉格朗日插值}\addcontentsline{toc}{subsection}{4.6 拉格朗日插值}

已知 $n+1$ 个点 $(x_i,y_i)$，求多项式在某点 $k$ 的值。

\[
f(k)=\sum_{i=0}^n y_i\prod_{j\ne i}\frac{k-x_j}{x_i-x_j}
\]

\begin{Verbatim}[fontsize=\small,frame=single,breaklines=true,breakanywhere=true]
ll Lagrange(vector<ll> x,vector<ll> y,ll k){
    int n=x.size()-1;
    ll ans=0;
    for(int i=0;i<=n;++i){
        ll num=1,den=1;
        for(int j=0;j<=n;++j) if(i!=j){
            num=num*(k-x[j])%p;
            den=den*(x[i]-x[j])%p;
        }
        ans=(ans+y[i]*num%p*qpow((den%p+p)%p,p-2,p))%p;
    }
    return (ans%p+p)%p;
}
\end{Verbatim}

\subsection*{4.7 Newton 迭代（多项式）}\addcontentsline{toc}{subsection}{4.7 Newton 迭代（多项式）}

常用于求根、求开方、求对数、求指数。

\[
g_{k+1}=g_k-\frac{f(g_k)}{f'(g_k)}
\]

\subsection*{4.8 生成函数常用模型}\addcontentsline{toc}{subsection}{4.8 生成函数常用模型}

\begin{itemize}[leftmargin=2em,itemsep=0.2em,topsep=0.2em]
\item 斐波那契：$\frac{x}{1-x-x^2}$
\item 等差级数：$\frac{x}{(1-x)^2}$
\item 幂和：$\sum_{n\ge 0}n^k x^n$ 可由算子法求出
\item 组合卷积：$\frac{1}{1-A(x)}=1+A(x)+A(x)^2+\cdots$


\end{itemize}
\section*{5. 数值算法}\addcontentsline{toc}{section}{5. 数值算法}

\subsection*{5.1 三分法}\addcontentsline{toc}{subsection}{5.1 三分法}

\begin{Verbatim}[fontsize=\small,frame=single,breaklines=true,breakanywhere=true]
double ternary(double l,double r){
    for(int it=0;it<200;++it){
        double m1=l+(r-l)/3,m2=r-(r-l)/3;
        if(f(m1)<f(m2)) l=m1;
        else r=m2;
    }
    return (l+r)/2;
}
\end{Verbatim}

\subsection*{5.2 Simpson 积分}\addcontentsline{toc}{subsection}{5.2 Simpson 积分}

\begin{Verbatim}[fontsize=\small,frame=single,breaklines=true,breakanywhere=true]
double simpson(double l,double r){
    double m=(l+r)/2;
    return (f(l)+4*f(m)+f(r))*(r-l)/6;
}

double asr(double l,double r,double eps,double whole){
    double m=(l+r)/2;
    double left=simpson(l,m),right=simpson(m,r);
    if(fabs(left+right-whole)<=15*eps) return left+right+(left+right-whole)/15;
    return asr(l,m,eps/2,left)+asr(m,r,eps/2,right);
}
\end{Verbatim}

\subsection*{5.3 牛顿迭代}\addcontentsline{toc}{subsection}{5.3 牛顿迭代}

\[
x_{k+1}=x_k-\frac{f(x_k)}{f'(x_k)}
\]


\subsection*{5.4 二分答案}\addcontentsline{toc}{subsection}{5.4 二分答案}

若答案满足单调性，则可二分枚举。

\begin{Verbatim}[fontsize=\small,frame=single,breaklines=true,breakanywhere=true]
bool check(double mid){
    return /* 单调判定 */;
}

double bs(double l,double r){
    for(int it=0;it<100;++it){
        double m=(l+r)/2;
        if(check(m)) r=m;
        else l=m;
    }
    return r;
}
\end{Verbatim}

\subsection*{5.5 牛顿迭代求平方根}\addcontentsline{toc}{subsection}{5.5 牛顿迭代求平方根}

\[
x_{k+1}=\frac{x_k+\frac{a}{x_k}}{2}
\]

\subsection*{5.6 黄金分割}\addcontentsline{toc}{subsection}{5.6 黄金分割}

\begin{Verbatim}[fontsize=\small,frame=single,breaklines=true,breakanywhere=true]
double golden(double l,double r){
    const double phi=(sqrt(5.0)-1)/2;
    double x1=r-phi*(r-l),x2=l+phi*(r-l);
    for(int it=0;it<200;++it){
        if(f(x1)<f(x2)) r=x2,x2=x1,x1=r-phi*(r-l);
        else l=x1,x1=x2,x2=l+phi*(r-l);
    }
    return (l+r)/2;
}
\end{Verbatim}


\section*{6. 博弈论}\addcontentsline{toc}{section}{6. 博弈论}

\subsection*{6.1 Nim 游戏}\addcontentsline{toc}{subsection}{6.1 Nim 游戏}

\begin{Verbatim}[fontsize=\small,frame=single,breaklines=true,breakanywhere=true]
int sg=0;
for(int x:a) sg^=x;
cout<<(sg?"Alice":"Bob")<<'\n';
\end{Verbatim}

\subsection*{6.2 SG 函数}\addcontentsline{toc}{subsection}{6.2 SG 函数}

\begin{Verbatim}[fontsize=\small,frame=single,breaklines=true,breakanywhere=true]
int dfs(int u){
    if(~sg[u]) return sg[u];
    set<int> s;
    for(int v:g[u]) s.insert(dfs(v));
    int mex=0;
    while(s.count(mex)) ++mex;
    return sg[u]=mex;
}
\end{Verbatim}

\subsection*{6.3 DAG 博弈}\addcontentsline{toc}{subsection}{6.3 DAG 博弈}

若状态转移图是 DAG，则可以按拓扑序/记忆化搜索求 SG 值。

\begin{Verbatim}[fontsize=\small,frame=single,breaklines=true,breakanywhere=true]
vector<int> g[N];
int sg[N];

int dfs(int u){
    if(~sg[u]) return sg[u];
    set<int> s;
    for(int v:g[u]) s.insert(dfs(v));
    int mex=0;
    while(s.count(mex)) ++mex;
    return sg[u]=mex;
}
\end{Verbatim}

\subsection*{6.4 威佐夫博弈}\addcontentsline{toc}{subsection}{6.4 威佐夫博弈}
\begin{Verbatim}[fontsize=\small,frame=single,breaklines=true,breakanywhere=true]
#include<bits/stdc++.h>
using namespace std;
const double l=(sqrt(5.0)+1.0)/2.0;
int n,m;
int main(){
    scanf("%d %d",&n,&m);
    if(n<m) swap(n,m);
    int d=n-m;
    if(int((double)d*l)==m) puts("0");
    else puts("1");
    return 0;
}
\end{Verbatim}

\subsection*{6.5 斐波那契尼姆博弈}\addcontentsline{toc}{subsection}{6.5 斐波那契尼姆博弈}
\begin{Verbatim}[fontsize=\small,frame=single,breaklines=true,breakanywhere=true]
#include<bits/stdc++.h>
long long n,x,y,z;
int main(){
    scanf("%lld",&n);
	while(66){
		if(n==1) return 0&puts("1");
		if(n==2) return 0&puts("2");
		x=1,y=2,z=3;
		while(z<n) x=y,y=z,z=x+y;
		if(z==n) return 0&printf("%lld",z); 
        else n-=y;
	}
	return 0;
}
\end{Verbatim}

\subsection*{6.6 mex 性质}\addcontentsline{toc}{subsection}{6.6 mex 性质}

\[
SG(u)=\mathrm{mex}\{SG(v)\mid u\to v\}
\]

其中 mex 表示最小非负整数。

\subsection*{6.7 记忆化模板}\addcontentsline{toc}{subsection}{6.7 记忆化模板}

\begin{Verbatim}[fontsize=\small,frame=single,breaklines=true,breakanywhere=true]
int win[N];
int solve(int u){
    if(win[u]!=-1) return win[u];
    for(int v:g[u])
        if(!solve(v)) return win[u]=1;
    return win[u]=0;
}
\end{Verbatim}
\section*{7. 计算几何}\addcontentsline{toc}{section}{7. 计算几何}

\subsection*{7.1 计算几何基础模板}\addcontentsline{toc}{subsection}{7.1 计算几何基础模板}

\begin{Verbatim}[fontsize=\small,frame=single,breaklines=true,breakanywhere=true]
// 1) 默认使用 long double，提升精度稳定性。
// 2) 所有判等都通过 dcmp 与 eps 完成，避免直接比较浮点。
#include <bits/stdc++.h>
using namespace std;

using ld = long double;
const ld EPS = 1e-12;

int dcmp(ld x){
    if(fabsl(x) < EPS) return 0;
    return x < 0 ? -1 : 1;
}

struct Point{
    ld x,y;
    Point(ld _x=0,ld _y=0): x(_x),y(_y) {}

    Point operator+(const Point& o) const { return Point(x+o.x,y+o.y); }
    Point operator-(const Point& o) const { return Point(x-o.x,y-o.y); }
    Point operator*(ld k) const { return Point(x*k,y*k); }
    Point operator / (ld k) const { return Point(x/k,y/k); }

    bool operator == (const Point& o) const {
        return dcmp(x-o.x)==0 && dcmp(y-o.y)==0;
    }

    bool operator < (const Point& o) const {
        if(dcmp(x-o.x)!=0) return x<o.x;
        return y<o.y;
    }
};

using Vec = Point;

// 叉积 a x b
ld cross(const Vec& a,const Vec& b){ return a.x*b.y-a.y*b.x; }
// 点积 a · b
ld dot(const Vec& a,const Vec& b){ return a.x*b.x+a.y*b.y; }
// 向量长度
ld norm(const Vec& a){ return sqrtl(dot(a,a)); }

// 方向判断：a->b 与 a->c 的转向
ld orient(const Point& a,const Point& b,const Point& c){
    return cross(b-a,c-a);
}

// 点 p 是否在线段 [a,b] 上（含端点）
bool onSegment(const Point& p,const Point& a,const Point& b){
    return dcmp(cross(p-a,p-b))==0 && dcmp(dot(p-a,p-b))<=0;
}

// 线段 [a,b] 与 [c,d] 是否相交（含端点/重叠）
bool segIntersect(const Point& a,const Point& b,const Point& c,const Point& d){
    ld c1 = orient(a,b,c),c2 = orient(a,b,d);
    ld c3 = orient(c,d,a),c4 = orient(c,d,b);

    if(dcmp(c1)*dcmp(c2)<0 && dcmp(c3)*dcmp(c4)<0) return true; // 严格相交
    if(onSegment(c,a,b) || onSegment(d,a,b) || onSegment(a,c,d) || onSegment(b,c,d)) return true;
    return false;
}

// 直线 AB 与 CD 的交点（需先保证不平行）
bool lineIntersection(const Point& a,const Point& b,const Point& c,const Point& d,Point& out){
    Vec u = b-a,v = d-c;
    ld den = cross(u,v);
    if(dcmp(den)==0) return false; // 平行或重合
    ld t = cross(c-a,v) / den;
    out = a+u*t;
    return true;
}

// 点到直线 AB 的距离
ld distPointLine(const Point& p,const Point& a,const Point& b){
    if(a==b) return norm(p-a);
    return fabsl(cross(b-a,p-a)) / norm(b-a);
}

// 点到线段 AB 的距离
ld distPointSegment(const Point& p,const Point& a,const Point& b){
    if(a==b) return norm(p-a);
    if(dcmp(dot(p-a,b-a))<0) return norm(p-a);
    if(dcmp(dot(p-b,a-b))<0) return norm(p-b);
    return distPointLine(p,a,b);
}

// 点 p 在直线 AB 上的投影点
Point projection(const Point& p,const Point& a,const Point& b){
    Vec v = b-a;
    if(a==b) return a;
    ld t = dot(p-a,v) / dot(v,v);
    return a+v*t;
}

// 点 p 关于直线 AB 的对称点
Point reflection(const Point& p,const Point& a,const Point& b){
    Point q = projection(p,a,b);
    return q*2-p;
}

// 多边形有向面积（逆时针为正）
ld polygonArea2(const vector<Point>& p){
    int n = (int)p.size();
    ld s = 0;
    for(int i=0;i<n;++i){
        s += cross(p[i],p[(i+1)%n]);
    }
    return s; // = 2*area
}

ld polygonArea(const vector<Point>& p){
    return fabsl(polygonArea2(p)) / 2.0;
}

// 点在多边形：0=外部,1=内部,2=边界（多边形可为非凸）
int pointInPolygon(const vector<Point>& poly,const Point& a){
    int n = (int)poly.size();
    bool in = false;
    for(int i=0,j=n-1;i<n;j=i++){
        const Point &p=poly[i],&q=poly[j];
        if(onSegment(a,p,q)) return 2;
        bool inter = ((dcmp(p.y-a.y)>0)!=(dcmp(q.y-a.y)>0)) &&
                     (dcmp(a.x-(q.x-p.x)*(a.y-p.y)/(q.y-p.y)-p.x) < 0);
        if(inter) in = !in;
    }
    return in ? 1 : 0;
}

// Andrew 凸包（返回逆时针，不重复首尾点）
vector<Point> convexHull(vector<Point> p){
    sort(p.begin(),p.end());
    p.erase(unique(p.begin(),p.end(),[](const Point& a,const Point& b){return a==b;}),p.end());
    int n = (int)p.size();
    if(n<=1) return p;

    vector<Point> st;
    // 下凸壳
    for(int i=0;i<n;++i){
        while((int)st.size()>=2 && dcmp(orient(st[(int)st.size()-2],st.back(),p[i]))<=0) st.pop_back();
        st.push_back(p[i]);
    }
    // 上凸壳
    int t = (int)st.size();
    for(int i=n-2;i>=0;--i){
        while((int)st.size()>t && dcmp(orient(st[(int)st.size()-2],st.back(),p[i]))<=0) st.pop_back();
        st.push_back(p[i]);
    }
    st.pop_back();
    return st;
}
\end{Verbatim}

\subsection*{7.2 凸包相关}\addcontentsline{toc}{subsection}{7.2 凸包相关}

\subsubsection*{7.2.1 旋转卡壳求凸包直径}\addcontentsline{toc}{subsubsection}{7.2.1 旋转卡壳求凸包直径}

\begin{Verbatim}[fontsize=\small,frame=single,breaklines=true,breakanywhere=true]
// 输入为逆时针凸包 ch（无重复首尾点）
// 返回最远点对的平方距离
long double convexDiameter2(const vector<Point>& ch){
    int n=(int)ch.size();
    if(n==1) return 0;
    if(n==2) return dot(ch[0]-ch[1],ch[0]-ch[1]);
    int j=1;
    long double ans=0;
    for(int i=0;i<n;++i){
        int ni=(i+1)%n;
        while(fabsl(cross(ch[ni]-ch[i],ch[(j+1)%n]-ch[i])) > fabsl(cross(ch[ni]-ch[i],ch[j]-ch[i])))
            j=(j+1)%n;
        ans=max(ans,dot(ch[i]-ch[j],ch[i]-ch[j]));
        ans=max(ans,dot(ch[ni]-ch[j],ch[ni]-ch[j]));
    }
    return ans;
}
\end{Verbatim}
\subsubsection*{7.2.2 Minkowski 和（凸多边形）}\addcontentsline{toc}{subsubsection}{7.2.2 Minkowski 和（凸多边形）}

\begin{Verbatim}[fontsize=\small,frame=single,breaklines=true,breakanywhere=true]
// 说明：输入 A/B 必须是逆时针凸包，且首元素为字典序最小点。
// 输出 C = A (+) B 的逆时针凸包（不重复首尾）。
vector<Point> minkowskiSum(vector<Point> A,vector<Point> B){
    auto shiftMin = [&](vector<Point>& P){
        int n=P.size(),id=0;
        for(int i=1;i<n;++i) if(P[i]<P[id]) id=i;
        rotate(P.begin(),P.begin()+id,P.end());
    };
    shiftMin(A); shiftMin(B);
    int n=A.size(),m=B.size();
    vector<Vec> EA,EB;
    for(int i=0;i<n;++i) EA.push_back(A[(i+1)%n]-A[i]);
    for(int i=0;i<m;++i) EB.push_back(B[(i+1)%m]-B[i]);

    vector<Point> C;
    C.push_back(A[0]+B[0]);
    int i=0,j=0;
    while(i<n || j<m){
        Vec va = (i<n?EA[i]:Vec(0,0));
        Vec vb = (j<m?EB[j]:Vec(0,0));
        if(i<n && (j==m || dcmp(cross(va,vb))>0)) C.push_back(C.back()+va),++i;
        else if(j<m && (i==n || dcmp(cross(va,vb))<0)) C.push_back(C.back()+vb),++j;
        else C.push_back(C.back()+va+vb),++i,++j;
    }
    C.pop_back();
    return convexHull(C); // 去共线冗余点
}
\end{Verbatim}

\subsection*{7.3 圆与直线}\addcontentsline{toc}{subsection}{7.3 圆与直线}

\begin{Verbatim}[fontsize=\small,frame=single,breaklines=true,breakanywhere=true]
struct Circle{
    Point o;
    ld r;
    Circle(Point _o=Point(),ld _r=0): o(_o),r(_r) {}
};

// 直线 AB 与圆 C 的交点个数（0/1/2），交点存入 out
int lineCircleIntersection(const Point& a,const Point& b,const Circle& C,vector<Point>& out){
    out.clear();
    Point h = projection(C.o,a,b);           // 圆心到直线投影
    ld d = norm(h-C.o);
    if(dcmp(d-C.r) > 0) return 0;
    Vec v = b-a;
    ld len = norm(v);
    if(dcmp(len)==0) return 0;
    v = v / len;
    if(dcmp(d-C.r)==0){
        out.push_back(h);
        return 1;
    }
    ld t = sqrtl(max((ld)0,C.r*C.r-d*d));
    out.push_back(h+v*t);
    out.push_back(h-v*t);
    return 2;
}

// 两圆位置关系：
// 0: 外离,1: 外切,2: 相交,3: 内切,4: 内含
int circleRelation(const Circle& A,const Circle& B){
    ld d = norm(A.o-B.o);
    if(dcmp(d-(A.r+B.r))>0) return 0;
    if(dcmp(d-(A.r+B.r))==0) return 1;
    if(dcmp(d-fabsl(A.r-B.r))>0) return 2;
    if(dcmp(d-fabsl(A.r-B.r))==0) return 3;
    return 4;
}
\end{Verbatim}

\subsection*{7.4 最近点对（分治）}\addcontentsline{toc}{subsection}{7.4 最近点对（分治）}

\begin{Verbatim}[fontsize=\small,frame=single,breaklines=true,breakanywhere=true]
// 最近点对 O(n log n)（返回最小距离平方）
long double closestPairRec(vector<Point>& p,int l,int r,vector<Point>& tmp){
    if(r-l<=3){
        long double ans=1e100;
        for(int i=l;i<=r;++i)
            for(int j=i+1;j<=r;++j)
                ans=min(ans,dot(p[i]-p[j],p[i]-p[j]));
        sort(p.begin()+l,p.begin()+r+1,[](const Point& a,const Point& b){return a.y<b.y;});
        return ans;
    }
    int m=(l+r)>>1;
    ld midx=p[m].x;
    ld d=min(closestPairRec(p,l,m,tmp),closestPairRec(p,m+1,r,tmp));
    merge(p.begin()+l,p.begin()+m+1,p.begin()+m+1,p.begin()+r+1,tmp.begin(),
          [](const Point& a,const Point& b){return a.y<b.y;});
    copy(tmp.begin(),tmp.begin()+(r-l+1),p.begin()+l);

    int ts=0;
    for(int i=l;i<=r;++i) if((p[i].x-midx)*(p[i].x-midx) < d) tmp[ts++]=p[i];
    for(int i=0;i<ts;++i)
        for(int j=i+1;j<ts && (tmp[j].y-tmp[i].y)*(tmp[j].y-tmp[i].y)<d;++j)
            d=min(d,dot(tmp[i]-tmp[j],tmp[i]-tmp[j]));
    return d;
}

long double closestPair2(vector<Point> p){
    sort(p.begin(),p.end(),[](const Point& a,const Point& b){
        if(dcmp(a.x-b.x)!=0) return a.x<b.x;
        return a.y<b.y;
    });
    vector<Point> tmp(p.size());
    return closestPairRec(p,0,(int)p.size()-1,tmp);
}
\end{Verbatim}


\section*{8.杂项}\addcontentsline{toc}{section}{8.杂项}

\subsection*{8.1 高精度全家桶}\addcontentsline{toc}{subsection}{8.1 高精度全家桶}

\begin{Verbatim}[fontsize=\small,frame=single,breaklines=true,breakanywhere=true]
// OI-wiki
struct BigIntTiny {
    int sign;
    std::vector<int> v;

    BigIntTiny() : sign(1) {}
    BigIntTiny(const std::string &s) { *this = s; }
    BigIntTiny(int v) {
        char buf[21];
        sprintf(buf,"%d",v);
        *this = buf;
    }
    void zip(int unzip) {
        if (unzip == 0) {
            for (int i = 0; i < (int)v.size(); i++)
                v[i] = get_pos(i*4)+get_pos(i*4+1)*10+get_pos(i*4+2)*100+get_pos(i*4+3)*1000;
        } else
            for (int i = (v.resize(v.size()*4),(int)v.size()-1),a; i >= 0; i--)
                a = (i % 4 >= 2) ? v[i / 4] / 100 : v[i / 4] % 100,v[i] = (i & 1) ? a / 10 : a % 10;
        setsign(1,1);
    }
    int get_pos(unsigned pos) const { return pos >= v.size() ? 0 : v[pos]; }
    BigIntTiny &setsign(int newsign,int rev) {
        for (int i = (int)v.size()-1; i > 0 && v[i] == 0; i--)
            v.erase(v.begin()+i);
        sign = (v.size() == 0 || (v.size() == 1 && v[0] == 0)) ? 1 : (rev ? newsign*sign : newsign);
        return *this;
    }
    std::string to_str() const {
        BigIntTiny b = *this;
        std::string s;
        for (int i = (b.zip(1),0); i < (int)b.v.size(); ++i)
            s += char(*(b.v.rbegin()+i)+'0');
        return (sign < 0 ? "-" : "")+(s.empty() ? std::string("0") : s);
    }
    bool absless(const BigIntTiny &b) const {
        if (v.size() != b.v.size()) return v.size() < b.v.size();
        for (int i = (int)v.size()-1; i >= 0; i--)
            if (v[i] != b.v[i]) return v[i] < b.v[i];
        return false;
    }
    BigIntTiny operator-() const {
        BigIntTiny c = *this;
        c.sign = (v.size() > 1 || v[0]) ? -c.sign : 1;
        return c;
    }
    BigIntTiny &operator=(const std::string &s) {
        if (s[0] == '-')
            *this = s.substr(1);
        else {
            for (int i = (v.clear(),0); i < (int)s.size(); ++i)
                v.push_back(*(s.rbegin()+i)-'0');
            zip(0);
        }
        return setsign(s[0] == '-' ? -1 : 1,sign = 1);
    }
    bool operator<(const BigIntTiny &b) const {
        return sign != b.sign ? sign < b.sign : (sign == 1 ? absless(b) : b.absless(*this));
    }
    bool operator==(const BigIntTiny &b) const { return v == b.v && sign == b.sign; }
    BigIntTiny &operator+=(const BigIntTiny &b) {
        if (sign != b.sign) return *this = (*this)--b;
        v.resize(std::max(v.size(),b.v.size())+1);
        for (int i = 0,carry = 0; i < (int)b.v.size() || carry; i++) {
            carry += v[i]+b.get_pos(i);
            v[i] = carry % 10000,carry /= 10000;
        }
        return setsign(sign,0);
    }
    BigIntTiny operator+(const BigIntTiny &b) const {
        BigIntTiny c = *this;
        return c += b;
    }
    void add_mul(const BigIntTiny &b,int mul) {
        v.resize(std::max(v.size(),b.v.size())+2);
        for (int i = 0,carry = 0; i < (int)b.v.size() || carry; i++) {
            carry += v[i]+b.get_pos(i)*mul;
            v[i] = carry % 10000,carry /= 10000;
        }
    }
    BigIntTiny operator-(const BigIntTiny &b) const {
        if (b.v.empty() || b.v.size() == 1 && b.v[0] == 0) return *this;
        if (sign != b.sign) return (*this)+-b;
        if (absless(b)) return -(b-*this);
        BigIntTiny c;
        for (int i = 0,borrow = 0; i < (int)v.size(); i++) {
            borrow += v[i]-b.get_pos(i);
            c.v.push_back(borrow);
            c.v.back() -= 10000*(borrow >>= 31);
        }
        return c.setsign(sign,0);
    }
    BigIntTiny operator*(const BigIntTiny &b) const {
        if (b < *this) return b**this;
        BigIntTiny c,d = b;
        for (int i = 0; i < (int)v.size(); i++,d.v.insert(d.v.begin(),0))
            c.add_mul(d,v[i]);
        return c.setsign(sign*b.sign,0);
    }
    BigIntTiny operator/(const BigIntTiny &b) const {
        BigIntTiny c,d;
        BigIntTiny e=b;
        e.sign=1;

        d.v.resize(v.size());
        double db = 1.0 / (b.v.back()+(b.get_pos((unsigned)b.v.size()-2) / 1e4) +
                           (b.get_pos((unsigned)b.v.size()-3)+1) / 1e8);
        for (int i = (int)v.size()-1; i >= 0; i--) {
            c.v.insert(c.v.begin(),v[i]);
            int m = (int)((c.get_pos((int)e.v.size())*10000+c.get_pos((int)e.v.size()-1))*db);
            c = c-e*m,c.setsign(c.sign,0),d.v[i] += m;
            while (!(c < e))
                c = c-e,d.v[i] += 1;
        }
        return d.setsign(sign*b.sign,0);
    }
    BigIntTiny operator%(const BigIntTiny &b) const { return *this-*this / b*b; }
    bool operator>(const BigIntTiny &b) const { return b < *this; }
    bool operator<=(const BigIntTiny &b) const { return !(b < *this); }
    bool operator>=(const BigIntTiny &b) const { return !(*this < b); }
    bool operator!=(const BigIntTiny &b) const { return !(*this == b); }
};
\end{Verbatim}

\subsection*{8.2快速平方根取倒数}\addcontentsline{toc}{subsection}{8.2快速平方根取倒数}

\begin{Verbatim}[fontsize=\small,frame=single,breaklines=true,breakanywhere=true]
loat Q_rsqrt( float number ){
    long i;
    float x2,y;
    const float threehalfs = 1.5F;
    x2 = number*0.5F;
    y  = number;
    i  =*( long*) &y;
    y  =*( float*) &i;
    return y;
}
\end{Verbatim}


\subsection*{8.3 模拟退火}\addcontentsline{toc}{subsection}{8.3 模拟退火}

\begin{Verbatim}[fontsize=\small,frame=single,breaklines=true,breakanywhere=true]
#include <cmath>
#include <cstdlib>
#include <ctime>
#include <iomanip>
#include <iostream>

constexpr int N = 10005;
int n, x[N], y[N], w[N];
double ansx, ansy, dis;

double Rand() { return (double)rand() / RAND_MAX; }

double calc(double xx, double yy) {
  double res = 0;
  for (int i = 1; i <= n; ++i) {
    double dx = x[i] - xx, dy = y[i] - yy;
    res += sqrt(dx * dx + dy * dy) * w[i];
  }
  if (res < dis) dis = res, ansx = xx, ansy = yy;
  return res;
}

void simulateAnneal() {
  double t = 100000;
  double nowx = ansx, nowy = ansy;
  while (t > 0.001) {
    double nxtx = nowx + t * (Rand() * 2 - 1);
    double nxty = nowy + t * (Rand() * 2 - 1);
    double delta = calc(nxtx, nxty) - calc(nowx, nowy);
    if (exp(-delta / t) > Rand()) nowx = nxtx, nowy = nxty;
    t *= 0.97;
  }
  for (int i = 1; i <= 1000; ++i) {
    double nxtx = ansx + t * (Rand() * 2 - 1);
    double nxty = ansy + t * (Rand() * 2 - 1);
    calc(nxtx, nxty);
  }
}

int main() {
  std::cin.tie(nullptr)->sync_with_stdio(false);
  srand(0);  // 注意，在实际使用中，不应使用固定的随机种子。
  std::cin >> n;
  for (int i = 1; i <= n; ++i) {
    std::cin >> x[i] >> y[i] >> w[i];
    ansx += x[i], ansy += y[i];
  }
  ansx /= n, ansy /= n, dis = calc(ansx, ansy);
  simulateAnneal();
  std::cout << std::fixed << std::setprecision(3) << ansx << ' ' << ansy
            << '\n';
  return 0;
}
\end{Verbatim}

\subsection*{8.4 快输快读（轻量化）}\addcontentsline{toc}{subsection}{8.4 快输快读（轻量化）}

\begin{Verbatim}[fontsize=\small,frame=single,breaklines=true,breakanywhere=true]
#include <bits/stdc++.h>
using namespace std;
template <class T> T read() {
    T f(1), x(0);
    int ch = getchar();
    while(ch < '0' || ch > '9') {
        if(ch == '-')
            f = -1;
        ch = getchar();
    }
    while(ch >= '0' && ch <= '9') {
        x = x * 10 + (ch ^ 48);
        ch = getchar();
    }
    return f * x;
}
template <class T> void write(T x) {
    if(x < 0)
        putchar('-'), x = -x;
    if(x > 9)
        write(x / 10);
    putchar(x % 10 + '0');
    return;
}
\end{Verbatim}

\subsection*{8.5 快输快读（重量级）}\addcontentsline{toc}{subsection}{8.5 快输快读（重量级）}

\begin{Verbatim}[fontsize=\small,frame=single,breaklines=true,breakanywhere=true]
#include <stdio.h>
#include <iostream>
const unsigned int bufl = 1 << 11;
const double base1[16] = {1,    1e-1, 1e-2,  1e-3,  1e-4,  1e-5,  1e-6,  1e-7,
                          1e-8, 1e-9, 1e-10, 1e-11, 1e-12, 1e-13, 1e-14, 1e-15};
const double base2[16] = {1,   1e1, 1e2,  1e3,  1e4,  1e5,  1e6,  1e7,
                          1e8, 1e9, 1e10, 1e11, 1e12, 1e13, 1e14, 1e15};
struct IN {
    FILE *IT;
    char ibuf[bufl], *is = ibuf, *it = ibuf;
    IN() {
        IT = stdin;
    }
    IN(char *a) {
        IT = fopen(a, "r");
    }
    inline char getChar() {
        if(is == it) {
            it = (is = ibuf) + fread(ibuf, 1, bufl, IT);
            if(is == it)
                return EOF;
        }
        return *is++;
    }
    template <typename Temp> inline void getInt(Temp &a) {
        a = 0;
        unsigned int b = 0, c = getChar();
        while(c < 48 || c > 57)
            b ^= (c == 45), c = getChar();
        while(c >= 48 && c <= 57)
            a = (a << 1) + (a << 3) + c - 48, c = getChar();
        if(b)
            a = -a;
    }
    IN &operator>>(char &a) {
        a = getChar();
        return *this;
    }
    IN &operator>>(char *a) {
        do {
            *a = getChar();
        } while(*a <= 32);
        while(*a > 32)
            *++a = getChar();
        *a = 0;
        return *this;
    }
    IN &operator>>(unsigned int &a) {
        getInt(a);
        return *this;
    }
};
struct OUT {
    FILE *IT;
    char obuf[bufl], *os = obuf, *ot = obuf + bufl;
    unsigned int Eps;
    long double Acc;
    OUT() {
        IT = stdout, Eps = 6, Acc = 1e-6;
    }
    OUT(char *a) {
        IT = fopen(a, "w"), Eps = 6, Acc = 1e-6;
    }
    inline void ChangEps(unsigned int x = 6) {
        Eps = x;
    }
    inline void flush() {
        fwrite(obuf, 1, os - obuf, IT);
        os = obuf;
    }
    inline void putChar(unsigned int a) {
        *os++ = a;
        if(os == ot)
            flush();
    }
    template <typename Temp> inline void putInt(Temp a) {
        if(a < 0) {
            putChar(45);
            a = -a;
        }
        if(a < 10) {
            putChar(a + 48);
            return;
        }
        putInt(a / 10);
        putChar(a % 10 + 48);
    }
    OUT &operator<<(char a) {
        putChar(a);
        return *this;
    }
    OUT &operator<<(char *a) {
        while(*a > 32)
            putChar(*a++);
        return *this;
    }
    OUT &operator<<(unsigned int a) {
        putInt(a);
        return *this;
    }
    ~OUT() {
        flush();
    }
};
IN fin;
OUT fout;
\end{Verbatim}

\section*{9. 数据结构}\addcontentsline{toc}{section}{9. 数据结构}

\subsection*{9.1 FHQ Treap}\addcontentsline{toc}{subsection}{9.1 FHQ Treap}

\begin{Verbatim}[fontsize=\small,frame=single,breaklines=true,breakanywhere=true]
struct FHQ {
    struct Node {
        int ls, rs, val, siz, rnd;
    }t[MAXN];
    int rt, cnt;
    void pushup(int p) {
        t[p].siz = t[t[p].ls].siz + t[t[p].rs].siz + 1;
    }
    int New(int a) {
        t[++cnt].val = a;
        t[cnt].siz = 1;
        t[cnt].rnd = Rand();
        return cnt;
    }
    void split(int now, int k, int &x, int &y) {
        if (!now) {
            x = y = 0;
            return;
        }
        if (t[now].val <= k) {
            x = now;
            split(t[now].rs, k, t[now].rs, y);
        }
        else {
            y = now;
            split(t[now].ls, k, x, t[now].ls);
        }
        pushup(now);
    }
    int merge(int u, int v) {
        if (!u || !v) return u | v;
        if (t[u].rnd < t[v].rnd) {
            t[u].rs = merge(t[u].rs, v);
            pushup(u);
            return u;
        }
        t[v].ls = merge(u, t[v].ls);
        pushup(v);
        return v;
    }
    void insert(int a) {
        int x, y;
        split(rt, a - 1, x, y);
        rt = merge(merge(x, New(a)), y);
    }
    void erase(int a) {
        int x, y, z;
        split(rt, a - 1, x, y);
        split(y, a, y, z);
        y = merge(t[y].ls, t[y].rs);
        rt = merge(merge(x, y), z);
    }
    int rk(int a) {
        int x, y;
        split(rt, a - 1, x, y);
        int ret = t[x].siz + 1;
        rt = merge(x, y);
        return ret;
    }
    int kth(int now, int k) {
        if (k <= t[t[now].ls].siz) return kth(t[now].ls, k);
        if (t[t[now].ls].siz + 1 < k) return kth(t[now].rs, k - t[t[now].ls].siz - 1);
        return t[now].val; 
    }
    int pre(int a) {
        int x, y;
        split(rt, a - 1, x, y);
        int ret = kth(x, t[x].siz);
        rt = merge(x, y);
        return ret;
    }
    int nxt(int a) {
        int x, y;
        split(rt, a, x, y);
        int ret = kth(y, 1);
        rt = merge(x, y);
        return ret;
    }
}t;
\end{Verbatim}

\subsection*{9.2 CDQ 分治}\addcontentsline{toc}{subsection}{9.2 CDQ 分治}

\begin{Verbatim}[fontsize=\small,frame=single,breaklines=true,breakanywhere=true]
auto cdq = [&](int l, int r, auto &self) -> void {
    if (l == r) return;
    int mid = (l + r) >> 1;
    self(l, mid, self); self(mid + 1, r, self);
    int p1 = l, p2 = mid + 1, cur = l;
    while (p1 <= mid && p2 <= r) {
        if (a[p1].b <= a[p2].b) {
            t.update(a[p1].c, a[p1].w);
            tmp[cur++] = a[p1++];
        }
        else {
            a[p2].ans += t.query(a[p2].c);
            tmp[cur++] = a[p2++];
        }
    }
    while (p1 <= mid) {
        t.update(a[p1].c, a[p1].w);
        tmp[cur++] = a[p1++];
    }
    while (p2 <= r) {
        a[p2].ans += t.query(a[p2].c);
        tmp[cur++] = a[p2++];
    }
    for (int i = l; i <= mid; i++) t.update(a[i].c, -a[i].w);
    for (int i = l; i <= r; i++) a[i] = tmp[i];
};
\end{Verbatim}

\subsection*{9.3 ST 表}\addcontentsline{toc}{subsection}{9.3 ST 表}

\begin{Verbatim}[fontsize=\small,frame=single,breaklines=true,breakanywhere=true]
struct ST {
    int st[LOGN][MAXN];
    void build(int n, int *a) {
        int t = std::__lg(n) + 1;
        for (int i = 1; i <= n; i++) st[0][i] = a[i];
        for (int i = 1; i <= t; i++) {
            for (int j = 1; j <= n; j++) {
                st[i][j] = std::max(st[i - 1][j], st[i - 1][j + (1 << (i - 1))]);
            }
        }
    }
    int query(int l, int r) {
        int t = std::__lg(r - l + 1);
        return std::max(st[t][l], st[t][r - (1 << t) + 1]);
    }
};
\end{Verbatim}

\subsection*{9.4 树链剖分}\addcontentsline{toc}{subsection}{9.4 树链剖分}

\begin{Verbatim}[fontsize=\small,frame=single,breaklines=true,breakanywhere=true]
auto dfs1 = [&](auto self, int u, int f) -> void {
    siz[u] = 1; fa[u] = f; dep[u] = dep[f] + 1;
    for (const auto &v : G[u]) {
        if (v == f) continue;
        self(self, v, u);
        siz[u] += siz[v];
        if (siz[v] > siz[son[u]]) son[u] = v; 
    }
};
auto dfs2 = [&](auto self, int u, int tp) -> void {
    top[u] = tp;
    if (son[u]) self(self, son[u], tp);
    for (const auto &v : G[u]) {
        if (v == son[u] || v == fa[u]) continue;
        self(self, v, v);
    }
};
auto lca = [&](int u, int v) {
    while (top[u] != top[v]) {
        if (dep[top[u]] < dep[top[v]]) std::swap(u, v);
        u = fa[top[u]];
    }
    if (dep[u] > dep[v]) return v;
    return u;
};
\end{Verbatim}

\subsection*{9.5 树上倍增 / RMQ LCA（原模板）}\addcontentsline{toc}{subsection}{9.5 树上倍增 / RMQ LCA（原模板）}

\begin{Verbatim}[fontsize=\small,frame=single,breaklines=true,breakanywhere=true]
struct St {
    int st[MAXN][25];
    int dfnmin(int x, int y) {
        if (dfn[x] < dfn[y]) return x;
        return y; 
    }
    void init() {
        int t = std::__lg(n) + 1;
        for (int i = 1; i <= n; i++) st[dfn[i]][0] = fa[i];
        for (int j = 1; j <= t; j++) {
            for (int i = 1; i + (1 << (j - 1))<= n; i++) {
                st[i][j] = dfnmin(st[i][j - 1], st[i + (1 << (j - 1))][j - 1]);
            }
        }
    }
    int query(int l, int r) {
        int t = std::__lg(r - l + 1);
        return dfnmin(st[l][t], st[r - (1 << t) + 1][t]);
    }
}st;

auto lca = [&](int x, int y) {
    if (x == y) return x;
    if (dfn[x] > dfn[y]) std::swap(x, y);
    return st.query(dfn[x] + 1, dfn[y]);
};
\end{Verbatim}


\subsection*{9.6 线段树优化建图}\addcontentsline{toc}{subsection}{9.6 线段树优化建图}

\begin{Verbatim}[fontsize=\small,frame=single,breaklines=true,breakanywhere=true]
#include <bits/stdc++.h>

using ll = long long;
using pii = std::pair <int, int>;

constexpr int MAXN = 7e5 + 5;
constexpr int K = 3e6 + 5;

std::vector <pii> G[K << 1];

void add(int u, int v, int w) {
    G[u].push_back({v, w});
}

int tid[MAXN];

struct Sgt {
    struct Node {
        int l, r;
    }t[MAXN << 2];
    void build(int p, int l, int r) {
        t[p].l = l; t[p].r = r;
        if (l == r) {
            tid[l] = p;
            return;
        }
        int mid = (l + r) >> 1;
        add(p, p << 1, 0); add(p, p << 1 | 1, 0);
        add((p << 1) + K, p + K, 0); add((p << 1 | 1) + K, p + K, 0);
        build(p << 1, l, mid); build(p << 1 | 1, mid + 1, r);
    }
    void connect(int p, int u, int l, int r, int w, int op) {//op=0 v->[l, r] op=1 [l,r]->v
        if (l <= t[p].l && t[p].r <= r) {
            if (op == 0) add(tid[u] + K, p, w);
            else if (op == 1) add(p + K, tid[u], w);
            return;
        }
        int mid = (t[p].l + t[p].r) >> 1;
        if (l <= mid) connect(p << 1, u, l, r, w, op);
        if (r > mid) connect(p << 1 | 1, u, l, r, w, op);
    }
}t;

void solve() {
    int n, m, s;
    std::cin >> n >> m >> s;
    t.build(1, 1, n + 2 * m);
    for (int i = 1; i <= n + 2 * m; i++) {
        add(tid[i], tid[i] + K, 0);
        add(tid[i] + K, tid[i], 0);
    }
    for (int i = 1; i <= m; i++) {
        int l1, r1, l2, r2;
        std::cin >> l1 >> r1 >> l2 >> r2;
        t.connect(1, n + 2 * i - 1, l1, r1, 1, 1);
        t.connect(1, n + 2 * i - 1, l2, r2, 0, 0);
        t.connect(1, n + 2 * i, l2, r2, 1, 1);
        t.connect(1, n + 2 * i, l1, r1, 0, 0);
        
    }
    std::vector <int> dis(K << 1);
    auto bfs = [&](int s) {
        std::fill(dis.begin(), dis.end(), 0x3f3f3f3f);
        std::deque <int> q;
        q.push_back(tid[s]); dis[tid[s]] = 0;
        while (q.size()) {
            int u = q.front(); q.pop_front();
            for (const auto &p : G[u]) {
                int v = p.first, w = p.second;
                if (dis[v] > dis[u] + w) {
                    dis[v] = dis[u] + w;
                    if (w == 0) q.push_front(v);
                    else q.push_back(v);
                } 
            }
        }
    };
    bfs(s);
    for (int i = 1; i <= n; i++) std::cout << dis[tid[i]] << '\n';
}

int main() {
    std::ios::sync_with_stdio(0);
    std::cin.tie(0);
    std::cout.tie(0);

    int T = 1;
    // std::cin >> T;
    while (T--) {
        solve();
    }

    return 0;
}
\end{Verbatim}

\subsection*{9.7 回滚莫队}\addcontentsline{toc}{subsection}{9.7 回滚莫队}

\begin{Verbatim}[fontsize=\small,frame=single,breaklines=true,breakanywhere=true]
#include<bits/stdc++.h>
using namespace std;

int len;

struct Query {
    int l, r, id;
    bool operator < (const Query &rhs) {
        if (l / len != rhs.l / len) return l / len < rhs.l / len;
        return r < rhs.r;
    }
};

void solve() {
    int n;
    std::cin >> n;
    std::vector <int> a(n + 1), lsh; 
    for (int i = 1; i <= n; i++) {
        std::cin >> a[i];
        lsh.push_back(a[i]);
    }
    std::sort(lsh.begin(), lsh.end());
    lsh.erase(std::unique(lsh.begin(), lsh.end()), lsh.end());
    for (int i = 1; i <= n; i++) {
        a[i] = std::lower_bound(lsh.begin(), lsh.end(), a[i]) - lsh.begin() + 1;
    }

    int m;
    std::cin >> m;
    std::vector <Query> qry;
    for (int i = 1; i <= m; i++) {
        int l, r;
        std::cin >> l >> r;
        qry.push_back({l, r, i});
    }
    len = sqrt(n);
    std::sort(qry.begin(), qry.end());
    std::vector <int> lpos(n + 1), rpos(n + 1), tmp(n + 1), ans(m + 1);
    int lst = -1;
    int l, r, curans;
    for (int i = 0; i < m; i++) {
        if (qry[i].l / len == qry[i].r / len) {
            for (int j = qry[i].l; j <= qry[i].r; j++) {
                if (!tmp[a[j]]) tmp[a[j]] = j;
                else ans[qry[i].id] = std::max(ans[qry[i].id], j - tmp[a[j]]);
            }
            for (int j = qry[i].l; j <= qry[i].r; j++) tmp[a[j]] = 0;
            continue;
        }
        int blk = qry[i].l / len;
        blk++;
        if (blk != lst) {
            r = len * blk - 1, l = r + 1;
            curans = 0;
            lst = blk;
            std::fill(lpos.begin(), lpos.end(), 0);
            std::fill(rpos.begin(), rpos.end(), 0);
        }
        while (r < qry[i].r) {
            ++r;
            if (!lpos[a[r]]) {
                lpos[a[r]] = rpos[a[r]] = r;
            }
            else {
                rpos[a[r]] = r;
                curans = std::max(curans, r - lpos[a[r]]);
            }
        }
        int tmp_ans = curans;
        int _l = l;
        std::vector <int> v;
        while (_l > qry[i].l) {
            --_l;
            if (!rpos[a[_l]]) {
                v.push_back(a[_l]);
                rpos[a[_l]] = _l;
            }
            else tmp_ans = std::max(tmp_ans, rpos[a[_l]] - _l); 
        }
        ans[qry[i].id] = tmp_ans;
        for (const auto &x : v) rpos[x] = 0;
    }

    for (int i = 1; i <= m; i++) {
        std::cout << ans[i] << '\n';
    }
}

\end{Verbatim}

\subsection*{9.8 带修莫队}\addcontentsline{toc}{subsection}{9.8 带修莫队}

\begin{Verbatim}[fontsize=\small,frame=single,breaklines=true,breakanywhere=true]
#include<bits/stdc++.h>
using namespace std;
using ll = long long;

int len;

struct Query {
    int l, r, id, tim;
    bool operator <(const Query &rhs) const {
        if (l / len != rhs.l / len) return l / len < rhs.l / len;
        if (r / len != rhs.r / len) return r / len < rhs.r / len;
        return tim < rhs.tim;
    }
};

struct Modify {
    int pos, val, tim;
};

void solve() {
    int n, m;
    std::cin >> n >> m;
    //qB = mn^2 / B^2 -> B = (mn^2/q)^ 1/3 
    std::vector <int> a(n + 1), ans(m + 1);
    for (int i = 1; i <= n; i++) std::cin >> a[i];
    int qcnt = 0, mcnt = 0;
    std::vector <Query> qry;
    std::vector <Modify> mdf;
    mdf.push_back({});
    for (int i = 1; i <= m; i++) {
        std::string s;
        int x, y;
        std::cin >> s >> x >> y;
        if (s == "Q") {
            ++qcnt;
            qry.push_back({x, y, qcnt, mcnt});
        }
        else {
            ++mcnt;
            mdf.push_back({x, y, mcnt});
        }
    }
    len = -1; 
    ll mincost = 1e18;
    for (int B = 1; B <= n; B++) {
        ll cost = 1ll * qcnt * B + 1ll * n * n / B + 1ll * mcnt * n * n / B / B;
        if (cost < mincost) {
            mincost = cost;
            len = B;
        }
    }
    std::sort(qry.begin(), qry.end());
    int l = 1, r = 0, t = 0;
    std::vector <int> cnt(1000005);
    int curans = 0;

    auto add = [&](int pos) {
        if (!cnt[a[pos]]++) curans++;
    };

    auto del = [&](int pos) {
        if (!--cnt[a[pos]]) curans--;
    };

    auto modify = [&](int tim) {
        int pos = mdf[tim].pos;
        if (l <= pos && pos <= r) del(pos);
        std::swap(a[pos], mdf[tim].val);
        if (l <= pos && pos <= r) add(pos);
    };

    for (const auto &p : qry) {
        int _l = p.l, _r = p.r, _t = p.tim;
        while (t < _t) modify(++t);
        while (_t < t) modify(t--);
        while (l < _l) del(l++);
        while (_l < l) add(--l);
        while (r < _r)  add(++r);
        while (_r < r) del(r--);
        ans[p.id] = curans;
    }

    for (int i = 1; i <= qcnt; i++) {
        std::cout << ans[i] << '\n';
    }
}

\end{Verbatim}

\subsection*{9.9 树状数组套权值线段树}\addcontentsline{toc}{subsection}{9.9 树状数组套权值线段树}

\begin{Verbatim}[fontsize=\small,frame=single,breaklines=true,breakanywhere=true]
#include<bits/stdc++.h>
using namespace std;

constexpr int MAXN = 200000 + 5;
constexpr int MAXV = 1000000000;

struct Sgt {
    struct Node {
        int ls, rs, cnt;
    }t[MAXN * 605];
    int tot = 0;
    void pushup(int p) {
        t[p].cnt = t[t[p].ls].cnt + t[t[p].rs].cnt;
    }
    void insert(int &p, int l, int r, int pos, int delta) {
        if (!p) p = ++tot;
        if (l == r) {
            t[p].cnt += delta;
            return;
        }
        int mid = (l + r) >> 1;
        if (pos <= mid) insert(t[p].ls, l, mid, pos, delta);
        else insert(t[p].rs, mid + 1, r, pos, delta);
        pushup(p);
    }
}t;

struct Bit {
    #define lowbit(x) ((x) & (-(x)))
    int rt[MAXN];
    void update(int pos, int v, int delta) {
        for (; pos < MAXN; pos += lowbit(pos)) t.insert(rt[pos], 0, MAXV, v, delta);
    }
    int query(int l, int r, int k) {
        l--;
        std::vector <int> vl, vr;
        for (; l; l -= lowbit(l)) vl.push_back(rt[l]);
        for (; r; r -= lowbit(r)) vr.push_back(rt[r]);
        int L = 0, R = MAXV;
        while (L != R) {
            int cnt = 0, mid = (L + R) >> 1;
            for (const auto &x : vr) {
                cnt += t.t[t.t[x].ls].cnt;
            }
            for (const auto &x : vl) {
                cnt -= t.t[t.t[x].ls].cnt;
            }
            if (k <= cnt) {
                R = mid;
                for (auto &x : vr) x = t.t[x].ls;
                for (auto &x : vl) x = t.t[x].ls;
            }
            else {
                L = mid + 1;
                k -= cnt;
                for (auto &x : vr) x = t.t[x].rs;
                for (auto &x : vl) x = t.t[x].rs;
            }
        }
        return L;
    }
}bit;

\end{Verbatim}

\subsection*{9.10 笛卡尔树}\addcontentsline{toc}{subsection}{9.10 笛卡尔树}

\begin{Verbatim}[fontsize=\small,frame=single,breaklines=true,breakanywhere=true]
#include<bits/stdc++.h>
using namespace std;
using ll = long long;

void solve() {
    int n;
    std::cin >> n;
    std::vector <int> a(n + 1);
    for (int i = 1; i <= n; i++) std::cin >> a[i];

    struct Node {
        int ls, rs;
    };
    std::vector <Node> t(n + 1);
    auto build = [&](int n, std::vector <int> &a) {
        std::stack <int> st;
        for (int i = 1; i <= n; i++) {
            int lst = 0;
            while (st.size() && a[i] < a[st.top()]) {
                lst = st.top(); 
                st.pop();
            } 
            if (st.size()) t[st.top()].rs = i;
            t[i].ls = lst;
            st.push(i);
        }
    };

    build(n, a);
    ll ansl = 0, ansr = 0;
    for (int i = 1; i <= n; i++) {
        ansl ^= (1ll * i * (t[i].ls + 1));
        ansr ^= (1ll * i * (t[i].rs + 1));
    }

    std::cout << ansl << " " << ansr << '\n';
}

\end{Verbatim}

\subsection*{9.11 李超树}\addcontentsline{toc}{subsection}{9.11 李超树}

\begin{Verbatim}[fontsize=\small,frame=single,breaklines=true,breakanywhere=true]
#include<bits/stdc++.h>
using namespace std;

constexpr int MAXN = 1e5 + 5;
constexpr int MAXV = 40000;

using ll = long long;
using pii = std::pair <int, int>;
#define y1 y114514
#define y2 y114515

struct Seg {
    int x0, x1, y0, y1;
    double calc(int x) {
        if (!(x0 <= x && x <= x1)) return 0.0;
        if (x0 == x1) return std::max(y0, y1);
        double k = 1.0 * (y1 - y0) / (x1 - x0), b = 1.0 * y0 - k * x0;
        return k * x + b;
    };
}s[MAXN];

constexpr double eps = 1e-9;
int cmp(double x, double y) {
    if (x - eps > y) return 1;
    if (x + eps < y) return -1;
    return 0;
}

struct Sgt {
    struct Node {
        int l, r;
        int sid;
    }t[MAXV << 2];
    void build(int p, int l, int r) {
        t[p].l = l; t[p].r = r;
        if (l == r) return;
        int mid = (l + r) >> 1;
        build(p << 1, l, mid); build(p << 1 | 1, mid + 1, r);
    }
    void update(int p, int sid) {
        if (t[p].sid == 0) {
            t[p].sid = sid;
            return;
        }
        int mid = (t[p].l + t[p].r) >> 1;
        int cmid = cmp(s[t[p].sid].calc(mid), s[sid].calc(mid));
        if (cmid == -1 || (cmid == 0 && t[p].sid > sid)) std::swap(t[p].sid, sid);
        int cl = cmp(s[t[p].sid].calc(t[p].l), s[sid].calc(t[p].l)), cr = cmp(s[t[p].sid].calc(t[p].r), s[sid].calc(t[p].r));
        if (cl == -1 || (cl == 0 && t[p].sid > sid)) update(p << 1, sid);
        if (cr == -1 || (cr == 0 && t[p].sid > sid)) update(p << 1 | 1, sid);
    }
    void insert(int p, int sid) {
        if (s[sid].x0 <= t[p].l && t[p].r <= s[sid].x1) {
            update(p, sid);
            return;
        } 
        int mid = (t[p].l + t[p].r) >> 1;
        if (s[sid].x0 <= mid) insert(p << 1, sid);
        if (s[sid].x1 > mid) insert(p << 1 | 1, sid);
    }
    int query(int p, int pos) {
        if (t[p].l == t[p].r) return t[p].sid;
        int mid = (t[p].l + t[p].r) >> 1, cur = t[p].sid;
        if (pos <= mid) {
            int lid = query(p << 1, pos), cl = cmp(s[lid].calc(pos), s[cur].calc(pos));
            if (cl == 1 || (cl == 0 && lid < cur)) cur = lid;
        }
        else {
            int rid = query(p << 1 | 1, pos), cr = cmp(s[rid].calc(pos), s[cur].calc(pos));
            if (cr == 1 || (cr == 0 && rid < cur)) cur = rid;
        }
        return cur;
    }
}t;

void solve() {
    t.build(1, 1, MAXV);
    int n;
    std::cin >> n;
    int lst = 0, cnt = 0;
    for (int i = 1; i <= n; i++) {
        int op;
        std::cin >> op;
        if (op == 0) {
            int x;
            std::cin >> x;
            x = (x + lst - 1) % 39989 + 1;
            std::cout << (lst = t.query(1, x)) << '\n';
        }
        else {
            ++cnt;
            std::cin >> s[cnt].x0 >> s[cnt].y0 >> s[cnt].x1 >> s[cnt].y1;
            s[cnt].x0 = (s[cnt].x0 + lst - 1) % 39989 + 1;
            s[cnt].x1 = (s[cnt].x1 + lst - 1) % 39989 + 1;
            s[cnt].y0 = (s[cnt].y0 + lst - 1) % 1000000000 + 1;
            s[cnt].y1 = (s[cnt].y1 + lst - 1) % 1000000000 + 1;
            if (s[cnt].x0 > s[cnt].x1) {
                std::swap(s[cnt].x0, s[cnt].x1);
                std::swap(s[cnt].y0, s[cnt].y1);
            }
            t.insert(1, cnt);
        }
    } 
}

\end{Verbatim}

\subsection*{9.12 线段树合并}\addcontentsline{toc}{subsection}{9.12 线段树合并}

\begin{Verbatim}[fontsize=\small,frame=single,breaklines=true,breakanywhere=true]
#include<bits/stdc++.h>
using namespace std;

struct Node {
    int l, r, sum;
};

vector<Node> tr;

void pushup(int p) {
    tr[p].sum = tr[tr[p].l].sum + tr[tr[p].r].sum;
}

int merge(int a, int b, int l, int r) {
  if (!a) return b;
  if (!b) return a;
  if (l == r) {
    tr[a].sum += tr[b].sum;
    return a;
  }
  int mid = (l + r) >> 1;
  tr[a].l = merge(tr[a].l, tr[b].l, l, mid);
  tr[a].r = merge(tr[a].r, tr[b].r, mid + 1, r);
  pushup(a);
  return a;
}

\end{Verbatim}

\section*{10. 图论}\addcontentsline{toc}{section}{10. 图论}

\subsection*{10.1 最短路（Dijkstra）}\addcontentsline{toc}{subsection}{10.1 最短路（Dijkstra）}

\begin{Verbatim}[fontsize=\small,frame=single,breaklines=true,breakanywhere=true]
auto dijkstra = [&](int s) {
    std::priority_queue <std::pair<int, int>, std::vector<std::pair<int, int> >, std::greater<std::pair<int, int> > > q;
    std::fill(dis.begin(), dis.end(), INF);
    dis[s] = 0;
    q.push({0, s});
    while (q.size()) {
        int u =  q.top().second; q.pop();
        if (vis[u]) continue;
        vis[u] = 1;
        for (const auto &p : G[u]) {
            int v = p.first, w = p.second;
            if (dis[v] > dis[u] + w) {
                dis[v] = dis[u] + w;
                q.push({dis[v], v});
            }
        }
    }
};
\end{Verbatim}

\subsection*{10.2 Dinic 最大流}\addcontentsline{toc}{subsection}{10.2 Dinic 最大流}

\begin{Verbatim}[fontsize=\small,frame=single,breaklines=true,breakanywhere=true]
struct Flow {
    int n;
    struct Edge {
        int to, w, rev;
    };
    std::vector <std::vector<Edge> > G;
    std::vector <int> dep, cur;
    Flow(int _n = 0) : n(_n) {
        G.resize(n + 1);
        dep.resize(n + 1);
        cur.resize(n + 1);
    }
    void add(int u, int v, int w) {
        int a = G[u].size(), b = G[v].size();
        G[u].push_back({v, w, b});
        G[v].push_back({u, 0, a});
    }
    bool bfs(int s, int t) {
        for (int i = 0; i <= n; i++) {
            dep[i] = 0;
            cur[i] = 0;
        }
        std::queue <int> q;
        q.push(s); dep[s] = 1;
        while(q.size()) {
            int u = q.front(); q.pop();
            for (const auto &[v, w, id] : G[u]) {
                if (dep[v] || !w) continue;
                dep[v] = dep[u] + 1;
                q.push(v);
            }
        }
        return (dep[t] != 0); 
    }
    static constexpr int INF = 0x7f7f7f7f;
    int dfs(int u, int t, int flow) {
        if (u == t) return flow;
        int rest = flow;
        for (int i = cur[u]; i < G[u].size(); i++) {
            const auto &[v, w, rev] = G[u][i];
            if (!w || dep[v] != dep[u] + 1) continue;
            int used = dfs(v, t, std::min(w, rest));
            rest -= used;
            G[u][i].w -= used; G[v][rev].w += used;
            cur[u] = i;
            if (!rest) return flow;  
        }
        return flow - rest;
    }
    ll dinic(int s, int t) {
        ll ret = 0;
        while (bfs(s, t)) {
            int flow;
            while ((flow = dfs(s, t, INF))) ret += flow;
        }
        return ret;
    }
};
\end{Verbatim}

\subsection*{10.3 费用流（MCMF）}\addcontentsline{toc}{subsection}{10.3 费用流（MCMF）}

\begin{Verbatim}[fontsize=\small,frame=single,breaklines=true,breakanywhere=true]
struct Flow {
    static constexpr int INF = 0x7f7f7f7f;
    int n;
    struct Edge {
        int to, w, rev, c;
    };
    std::vector <std::vector<Edge> > G;
    std::vector <int> dep, cur, dis, inq;
    Flow(int _n = 0) : n(_n) {
        G.resize(n + 1);
        dep.resize(n + 1);
        cur.resize(n + 1);
        dis.resize(n + 1);
        inq.resize(n + 1);
    }
    void clear(int _n) {
        for (int i = 0; i <= n; i++) G[i].clear();
        G.clear(); G.resize(_n + 1);
        dep.clear(); dep.resize(_n + 1, 0);
        cur.clear(); cur.resize(_n + 1, 0);
        dis.clear(); dis.resize(_n + 1, 0);
        inq.clear(); inq.resize(_n + 1, 0);
        n = _n;
    }
    void add(int u, int v, int w, int c) {
        int a = G[u].size(), b = G[v].size();
        G[u].push_back({v, w, b, c});
        G[v].push_back({u, 0, a, -c});
    }
    bool SPFA(int s, int t) {
        for (int i = 0; i <= n; i++) {
            dep[i] = 0;
            cur[i] = 0;
            inq[i] = 0;
            dis[i] = INF;
        }
        std::queue <int> q;
        q.push(s); dep[s] = 1; inq[s] = 1; dis[s] = 0;
        while(q.size()) {
            int u = q.front(); q.pop(); inq[u] = 0;
            for (const auto &[v, w, id, c] : G[u]) {
                if (w && dis[v] > dis[u] + c) {
                    dis[v] = dis[u] + c;
                    dep[v] = dep[u] + 1;
                    if (!inq[v]) {
                        q.push(v);
                        inq[v] = 1; 
                    }    
                }
            }
        }
        return (dep[t] != 0); 
    }
    ll mincost = 0;
    int dfs(int u, int t, int flow) {
        if (u == t) return flow;
        int rest = flow;
        for (int i = cur[u]; i < G[u].size(); i++) {
            const auto &[v, w, rev, c] = G[u][i];
            if (!w || dep[v] != dep[u] + 1 || dis[v] != dis[u] + c) continue;
            int used = dfs(v, t, std::min(w, rest));
            rest -= used;
            G[u][i].w -= used; G[v][rev].w += used;
            mincost += 1ll * c * used;
            cur[u] = i;
            if (!rest) return flow;  
        }
        return flow - rest;
    }
    std::pair<ll, ll> MCMF(int s, int t) {
        ll ret = 0;
        while (SPFA(s, t)) {
            int flow;
            while ((flow = dfs(s, t, INF))) ret += flow;
        }
        return std::make_pair(ret, mincost);
    }
};
\end{Verbatim}

\subsection*{10.4 Tarjan 强连通分量}\addcontentsline{toc}{subsection}{10.4 Tarjan 强连通分量}

\begin{Verbatim}[fontsize=\small,frame=single,breaklines=true,breakanywhere=true]
auto tarjan = [&](int u, auto self) -> void {
    dfn[u] = low[u] = ++dcnt;
    s.push(u); ins[u] = 1;
    for (const auto &v : G[u]) {
        if (!dfn[v]) {
            self(v, self);
            low[u] = std::min(low[u], low[v]);
        }
        else if (ins[v]) low[u] = std::min(low[u], dfn[v]);
    }
    if (dfn[u] == low[u]) {
        ++scc;
        int cur;
        do {
            cur = s.top(); s.pop(); ins[cur] = 0;
            belong[cur] = scc;
        }while (cur != u);
    }
};
\end{Verbatim}
\subsection*{10.5 Borvuka}\addcontentsline{toc}{subsection}{10.5 Borvuka}

\begin{Verbatim}[fontsize=\small,frame=single,breaklines=true,breakanywhere=true]
#include <bits/stdc++.h>

const int MAXN = 5005;
const int MAXM = 2e5 + 5;

struct Dsu {
    int fa[MAXN];
    void init(int n) {
        for (int i = 1; i <= n; i++) fa[i] = i;
    }
    int find(int x) {
        return fa[x] == x ? x : fa[x] = find(fa[x]);
    }
}dsu;

int n, m;

struct Edge {
    int u, v, w;
}E[MAXM];

int main() {
    std::ios::sync_with_stdio(0);
    std::cin.tie(0);
    std::cout.tie(0);

    std::cin >> n >> m;
    for (int i = 1;i <= m; i++) {
        std::cin >> E[i].u >> E[i].v >> E[i].w;
    }

    #define ll long long
    auto boruvka = [&]() {
        ll ans = 0;
        dsu.init(n);
        std::vector <bool> vis(m + 1);
        int ecnt = 0;
        while (1) {
            bool flag = 0;
            std::vector <int> mn(n + 1, 0x3f3f3f3f), id(n + 1, 0);
            for (int i = 1; i <= m; i++) {
                if (vis[i]) continue;
                int u = E[i].u, v = E[i].v, w = E[i].w;
                int x = dsu.find(u), y = dsu.find(v);
                if (x == y) continue;
                flag = 1;
                if (w < mn[x]) {
                    mn[x] = w;
                    id[x] = i;
                }
                if (w < mn[y]) {
                    mn[y] = w;
                    id[y] = i;
                }
            }
            if (!flag || ecnt == n - 1) break;
            for (int i = 1; i <= n; i++) {
                int x = dsu.find(i);
                if (!id[x] || vis[id[x]]) continue;
                int u = E[id[x]].u, v = E[id[x]].v;
                int X = dsu.find(u), Y = dsu.find(v);
                if (X == Y) continue;
                ++ecnt;
                ans += E[id[x]].w;
                dsu.fa[X] = Y;
                vis[id[x]] = 1;
            }
        }
        if (ecnt < n - 1) return -1ll;
        return ans;
    };

    ll ret = boruvka();
    if (ret == -1) std::cout << "orz" << '\n';
    else std::cout << boruvka() << '\n';
}
\end{Verbatim}

\subsection*{10.6 2sat}\addcontentsline{toc}{subsection}{10.6 2sat}

\begin{Verbatim}[fontsize=\small,frame=single,breaklines=true,breakanywhere=true]
#include <bits/stdc++.h>

const int MAXN = 1e6 + 5;

int n, m;
std::vector <int> G[MAXN * 2];

int main() {
    std::cin >> n >> m;
    for (int i = 1; i <= m; i++) {
        int a, b, xa, xb;
        std::cin >> a >> xa >> b >> xb;

        auto calc = [&](int i, int a) {
            if (a) return i + n;
            return i;
        };

        G[calc(a, xa ^ 1)].push_back(calc(b, xb));
        G[calc(b, xb ^ 1)].push_back(calc(a, xa));
    }

    int dcnt = 0, scc = 0;
    std::vector <int> dfn(2 * n + 1), low(2 * n + 1), belong(2 * n + 1);
    std::vector <bool> ins(2 * n + 1);
    std::stack <int> s;
    auto tarjan = [&](int u, auto self) -> void {
        dfn[u] = low[u] = ++dcnt;
        s.push(u); ins[u] = 1;
        for (const auto &v : G[u]) {
            if (!dfn[v]) {
                self(v, self);
                low[u] = std::min(low[u], low[v]);
            }
            else if (ins[v]) low[u] = std::min(low[u], dfn[v]);
        }
        if (dfn[u] == low[u]) {
            ++scc;
            int cur;
            do {
                cur = s.top(); s.pop(); ins[cur] = 0;
                belong[cur] = scc;
            }while (cur != u);
        }
    };
    for (int i = 1; i <= 2 * n; i++) {
        if (!dfn[i]) tarjan(i, tarjan);
    }
    for (int i = 1; i <= n; i++) {
        if (belong[i] == belong[i + n]) {
            std::cout << "IMPOSSIBLE" << '\n';
            return 0;
        } 
    }

    std::cout << "POSSIBLE" << '\n';
    for (int i = 1; i <= n; i++) {
        if (belong[i] > belong[i + n]) std::cout << 1 << " ";
        else std::cout << 0 << " ";
    }
    std::cout << '\n';
}
\end{Verbatim}

\section*{11. 字符串}\addcontentsline{toc}{section}{11. 字符串}

\subsection*{11.1 字符串哈希}\addcontentsline{toc}{subsection}{11.1 字符串哈希}

\begin{Verbatim}[fontsize=\small,frame=single,breaklines=true,breakanywhere=true]
#include<bits/stdc++.h>
using ull = unsigned long long;
ull base = 131;
ull mod1 = 212370440130137957, mod2 = 1e9 + 7;

ull get_hash1(std::string s) {
  int len = s.size();
  ull ans = 0;
  for (int i = 0; i < len; i++) ans = (ans * base + (ull)s[i]) % mod1;
  return ans;
}

ull get_hash2(std::string s) {
  int len = s.size();
  ull ans = 0;
  for (int i = 0; i < len; i++) ans = (ans * base + (ull)s[i]) % mod2;
  return ans;
}

bool cmp(const std::string s, const std::string t) {
  bool f1 = get_hash1(s) != get_hash1(t);
  bool f2 = get_hash2(s) != get_hash2(t);
  return f1 || f2;
}

// 双哈希
\end{Verbatim}

\subsection*{11.2 KMP 与 exKMP}\addcontentsline{toc}{subsection}{11.2 KMP 与 exKMP}

\begin{Verbatim}[fontsize=\small,frame=single,breaklines=true,breakanywhere=true]
#include<bits/stdc++.h>
using namespace std;
vector<int> prefix_function(string s) {
  int n = (int)s.length();
  vector<int> pi(n);
  for (int i = 1; i < n; i++) {
    int j = pi[i - 1];
    while (j > 0 && s[i] != s[j]) j = pi[j - 1];
    if (s[i] == s[j]) j++;
    pi[i] = j;
  }
  return pi;
}
vector<int> z_function(string s) {
  int n = (int)s.length();
  vector<int> z(n);
  for (int i = 1, l = 0, r = 0; i < n; ++i) {
    if (i <= r && z[i - l] < r - i + 1) {
      z[i] = z[i - l];
    } else {
      z[i] = max(0, r - i + 1);
      while (i + z[i] < n && s[z[i]] == s[i + z[i]]) ++z[i];
    }
    if (i + z[i] - 1 > r) l = i, r = i + z[i] - 1;
  }
  return z;
}
\end{Verbatim}

\subsection*{11.3 Manacher}\addcontentsline{toc}{subsection}{11.3 Manacher}

\begin{Verbatim}[fontsize=\small,frame=single,breaklines=true,breakanywhere=true]
#include<bits/stdc++.h>
using namespace std;
int main(){
    int n;
    vector<int> d1(n);
    string s;
    cin>>n>>s;
    for (int i = 0, l = 0, r = -1; i < n; i++) {
        int k = (i > r) ? 1 : min(d1[l + r - i], r - i + 1);
        while (0 <= i - k && i + k < n && s[i - k] == s[i + k]) {
            k++;
        }
        d1[i] = k--;
        if (i + k > r) {
            l = i - k;
            r = i + k;
        }
    }
    vector<int> d2(n);
    for (int i = 0, l = 0, r = -1; i < n; i++) {
        int k = (i > r) ? 0 : min(d2[l + r - i + 1], r - i + 1);
        while (0 <= i - k - 1 && i + k < n && s[i - k - 1] == s[i + k]) {
            k++;
        }
        d2[i] = k--;
        if (i + k > r) {
            l = i - k - 1;
            r = i + k;
        }
    }
}
\end{Verbatim}

\subsection*{11.4 后缀自动机}\addcontentsline{toc}{subsection}{11.4 后缀自动机}

\begin{Verbatim}[fontsize=\small,frame=single,breaklines=true,breakanywhere=true]
#include<bits/stdc++.h>
using namespace std;
constexpr int MAXLEN = 100000;
struct state {
  int len, link;
  std::map<char, int> next;
};
state st[MAXLEN * 2];
int sz, last;
void sam_init() {
  st[0].len = 0;
  st[0].link = -1;
  sz++;
  last = 0;
}
void sam_extend(char c) {
  int cur = sz++;
  st[cur].len = st[last].len + 1;
  int p = last;
  while (p != -1 && !st[p].next.count(c)) {
    st[p].next[c] = cur;
    p = st[p].link;
  }
  if (p == -1) {
    st[cur].link = 0;
  } else {
    int q = st[p].next[c];
    if (st[p].len + 1 == st[q].len) {
      st[cur].link = q;
    } else {
      int clone = sz++;
      st[clone].len = st[p].len + 1;
      st[clone].next = st[q].next;
      st[clone].link = st[q].link;
      while (p != -1 && st[p].next[c] == q) {
        st[p].next[c] = clone;
        p = st[p].link;
      }
      st[q].link = st[cur].link = clone;
    }
  }
  last = cur;
}
\end{Verbatim}

\end{document}

